\documentclass{article}

% if you need to pass options to natbib, use, e.g.:
%     \PassOptionsToPackage{numbers, compress}{natbib}
% before loading neurips_2026

% The authors should use one of these tracks.
% Before accepting by the NeurIPS conference, select one of the options below.
% 0. "default" for submission

% \usepackage{neurips_2026}
\usepackage[preprint]{neurips_2026}
% \usepackage[main, final]{neurips_2026}

% the "default" option is equal to the "main" option, which is used for the Main Track with double-blind reviewing.
% 1. "main" option is used for the Main Track
%  \usepackage[main]{neurips_2026}
% 2. "position" option is used for the Position Paper Track
%  \usepackage[position]{neurips_2026}
% 3. "eandd" option is used for the Evaluations & Datasets Track
 % \usepackage[eandd]{neurips_2026}
 % if you need to opt-in for a single-blind submission in the E&D track:
 %\usepackage[eandd, nonanonymous]{neurips_2026}
% 4. "creativeai" option is used for the Creative AI Track
%  \usepackage[creativeai]{neurips_2026}
% 5. "sglblindworkshop" option is used for the Workshop with single-blind reviewing
 % \usepackage[sglblindworkshop]{neurips_2026}
% 6. "dblblindworkshop" option is used for the Workshop with double-blind reviewing
%  \usepackage[dblblindworkshop]{neurips_2026}

% After being accepted, the authors should add "final" behind the track to compile a camera-ready version.
% 1. Main Track
 % \usepackage[main, final]{neurips_2026}
% 2. Position Paper Track
%  \usepackage[position, final]{neurips_2026}
% 3. Evaluations & Datasets Track
 % \usepackage[eandd, final]{neurips_2026}
% 4. Creative AI Track
%  \usepackage[creativeai, final]{neurips_2026}
% 5. Workshop with single-blind reviewing
%  \usepackage[sglblindworkshop, final]{neurips_2026}
% 6. Workshop with double-blind reviewing
%  \usepackage[dblblindworkshop, final]{neurips_2026}
% Note. For the workshop paper template, both \title{} and \workshoptitle{} are required, with the former indicating the paper title shown in the title and the latter indicating the workshop title displayed in the footnote.
% For workshops (5., 6.), the authors should add the name of the workshop, "\workshoptitle" command is used to set the workshop title.
% \workshoptitle{WORKSHOP TITLE}

% "preprint" option is used for arXiv or other preprint submissions
 % \usepackage[preprint]{neurips_2026}

% to avoid loading the natbib package, add option nonatbib:
%    \usepackage[nonatbib]{neurips_2026}

\usepackage[utf8]{inputenc} % allow utf-8 input
\usepackage[T1]{fontenc}    % use 8-bit T1 fonts
\usepackage{hyperref}       % hyperlinks
% --- autoref name customization (Option A) ---
\renewcommand{\equationautorefname}{Eq.}
\usepackage{url}            % simple URL typesetting
\usepackage{booktabs}       % professional-quality tables
\usepackage{graphicx}       % resize tables and figures
\usepackage{placeins}       % keep floats within the intended section
\usepackage{amsmath}        % math environments and symbols
\usepackage{amsfonts}       % blackboard math symbols
\usepackage{nicefrac}       % compact symbols for 1/2, etc.
\usepackage{microtype}      % microtypography
\usepackage[table, dvipsnames]{xcolor}         % colors

\usepackage{tabularx}
\usepackage{booktabs}
\usepackage{array}
\usepackage{adjustbox}   %压缩表格


\definecolor{myblue}{HTML}{0064E0}
\definecolor{myred}{HTML}{DC6F7B}
\usepackage{wrapfig} 
\hypersetup{
    colorlinks,
    linkcolor=myblue,
    citecolor=myblue,
    filecolor=Mulberry,
    urlcolor=RoyalBlue,
}

% Note. For the workshop paper template, both \title{} and \workshoptitle{} are required, with the former indicating the paper title shown in the title and the latter indicating the workshop title displayed in the footnote. 
\title{Benign in Isolation, Harmful in Composition: Security Risks in Agent Skill Ecosystems}


% The \author macro works with any number of authors. There are two commands
% used to separate the names and addresses of multiple authors: \And and \AND.
%
% Using \And between authors leaves it to LaTeX to determine where to break the
% lines. Using \AND forces a line break at that point. So, if LaTeX puts 3 of 4
% authors names on the first line, and the last on the second line, try using
% \AND instead of \And before the third author name.


% \author{Anonymous Author(s)}
% \author{
%   Yi Xie\textsuperscript{1} \\
%   East China Normal University, Shanghai, China
%   \And
%   Jiawei Du\textsuperscript{2*} \\
%   Centre for Frontier AI Research, A*STAR, Singapore
%   \AND
%   Yu Cheng\textsuperscript{1,3} \\
%   East China Normal University, Shanghai, China \\
%   Shanghai Innovation Institute, Shanghai, China
%   \And
%   Jiuan Zhou\textsuperscript{1} \\
%   East China Normal University, Shanghai, China
%   \And
%   Zhaoxia Yin\textsuperscript{1,*} \\
%   East China Normal University, Shanghai, China \\
%   \textsuperscript{*}Corresponding author.
% }
\author{
    \textbf{Yi Xie}\textsuperscript{1}, 
    \textbf{Jiawei Du}\textsuperscript{*,2}, 
    \textbf{Yu Cheng}\textsuperscript{1,3}, 
    \textbf{Jiuan Zhou}\textsuperscript{1}, 
    \textbf{Zhaoxia Yin}\textsuperscript{*,1} \\ \textsuperscript{1}East China Normal University, Shanghai, China \\ 
    \textsuperscript{2}Centre for Frontier AI Research A*STAR, Singapore \\ 
    \textsuperscript{3}Shanghai Innovation Institute, Shanghai, China \\ 
    \small\textsuperscript{*}Corresponding authors.
}


\begin{document}


\maketitle


\begin{abstract}
Skills are becoming the capability layer through which LLM agents turn plans into actions, but their use also introduces security risks, including data leakage, unauthorized operations, and tool misuse. Existing skill vetting typically reviews each skill in isolation, whereas real agent tasks often involve multiple skill invocations across a shared execution context. A risk therefore emerges: one skill’s output, trust signal, authorization cue, or side effect can be carried into a later skill invocation. We define this resulting gap as \textbf{SCR} (\textbf{S}kill \textbf{C}omposition \textbf{R}isk): skills that appear benign under isolated evaluation can become harmful along an activated composition path. \textbf{ SCR-Bench} (\textbf{S}kill \textbf{C}omposition \textbf{R}isk Bench) is introduced to evaluate this risk in controlled, sandboxed skill environments. Instead of relying on the textual intent or surface-level behavior of individual skills, SCR-Bench records downstream state changes and path-level outcomes across activated skill paths. SCR-Bench contains three sub-benchmarks: \textbf{SCR-CapFlow} for capability-flow composition, \textbf{SCR-TrustLift} for trust-transfer composition, and \textbf{SCR-AuthBlur} for authorization-confusion composition. Across SCR-Bench, composed skill paths expose risks that are largely absent under isolated evaluation: in SCR-CapFlow, ASR reaches 33.6\% under composition, compared with near-zero isolated baselines; in SCR-TrustLift, ASR reaches near saturation above 96.5\% on four of five backends; and in SCR-AuthBlur, the risky-approval rate increases by 71.8\% relative to the L0 isolated baseline under the L1 context setting. The experimental results reveal the core pattern of “Benign in Isolation, Harmful in Composition,” showing that agent skill security should be assessed at the level of activated paths rather than isolated artifacts. By introducing SCR and SCR-Bench, this work aims to support path-aware risk evaluation and defense in LLM agent skill ecosystems. Our benchmark is available at \url{https://anonymous.4open.science/r/SCR-Bench-B20F}.
\end{abstract}


\begin{figure}[t]
  \centering
  \includegraphics[width=\linewidth]{figures/intro-0612.pdf}
  \caption{\textbf{Path-level risk in SCR.} The figure shows three skills that are safe under isolated review: a file-audit skill, a security--review skill, and an access--manager skill. When the agent composes them under task context, their outputs can activate capability flow, trust transfer, or authorization confusion, exposing sensitive resources even though no single skill is harmful in isolation.}

  \label{fig:composition-risk-overview}
\end{figure}



\section{Introduction}

% Skills are increasingly becoming the action layer of LLM agents, enabling them to move beyond passive reasoning and orchestrate long-horizon workflows through reusable modules that interact with external services, software environments, and local state~\citep{jiang2026sokagenticskills}.

LLM agents extend LLM reasoning with external actions, including tool calls, API invocations, memory access, and reusable procedural modules~\citep{yao2023react,schick2023toolformer,qin2024toollm,patil2024gorilla}. Recent red-teaming work further shows that agent memory and retrieval components can themselves become attack surfaces, as poisoned long-term memory or knowledge bases may steer later planning and execution~\citep{chen2024agentpoison}. Building on this tool-use paradigm, agentic skills serve as reusable action units that encapsulate procedural knowledge, execution policies, and interfaces for long-horizon workflows~\citep{jiang2026sokagenticskills}. Yet these modules are often installed and executed under implicit trust with limited vetting, creating a governance gap: a recent large-scale study of real-world agent skills found that 26.1\% contain at least one vulnerability, including data-exfiltration and privilege-escalation patterns~\citep{liu2026donotmentionmaliciousskills}. As a result, ensuring that agents can safely use skills has become a central challenge for skill-enabled systems, directly constraining their reliability, deployability, and practical utility~\citep{bhardwaj2026formalanalysis}.

% Existing work has largely addressed skill-related risks through artifact-level security review, treating each skill or tool as the primary object of inspection. In the broader agent-tool setting, AgentDojo~\citep{yuan2026aegis} and Agent Security Bench~\citep{wang2026systematicsecurityopenclaw} show that agents can be hijacked through untrusted context, prompt injection, and adversarial tool use. Skill-specific studies further expose risks embedded in the skill artifact itself: Skill-Inject~\citep{schmotz2026skillinject} studies malicious instructions hidden in skill files, while HarmfulSkillBench~\citep{jiang2026harmfulskillbench} shows that overtly harmful skills can reduce refusal behavior and enable unsafe tasks such as data exfiltration or destructive actions. Recent defenses therefore audit skills before deployment or execution. SkillScan~\citep{liu2026agentskillswild} and SkillSieve~\citep{hou2026skillsieve} combine static analysis with LLM-based semantic triage to detect vulnerabilities in skill source code, metadata, and declared behavior, while SkillFortify~\citep{bhardwaj2026formalanalysis} uses abstract interpretation and capability-based confinement to verify that a skill cannot exceed its declared authority under explicit security assumptions. Together, these studies establish standalone artifact vetting as a necessary security primitive for detecting excessive authority requests, embedded malicious instructions or code, and misrepresented behavior.

% Existing work has largely treated each skill or tool as the primary unit of security review. In the broader agent-tool setting, AgentDojo~\citep{yuan2026aegis} and Agent Security Bench~\citep{zhang2025agent} show that agents can be hijacked through untrusted context, prompt injection, and adversarial tool use, while skill-specific studies such as Skill-Inject~\citep{schmotz2026skillinject} and HarmfulSkillBench~\citep{jiang2026harmfulskillbench} expose malicious instructions or overtly harmful functionality embedded in skill artifacts. Recent defenses therefore audit skills before deployment or execution: SkillScan~\citep{liu2026agentskillswild} and SkillSieve~\citep{hou2026skillsieve} combine static analysis with LLM-based semantic triage, and SkillFortify~\citep{bhardwaj2026formalanalysis} applies abstract interpretation and capability-based confinement to check whether a skill can exceed its declared authority. Together, these studies establish standalone artifact vetting as a necessary security primitive for detecting excessive authority requests, embedded malicious instructions or code, and misrepresented behavior.

Existing work has largely treated each skill or tool as the primary unit of security review. In the broader agent-tool setting, prompt-injection benchmarks have formalized how untrusted inputs can manipulate LLM-integrated applications~\citep{liu2024formalizing}, while InjecAgent studies indirect prompt injection in tool-integrated LLM agents~\citep{zhan2024injecagent}. AgentDojo~\citep{debenedetti2024agentdojo} and Agent Security Bench~\citep{zhang2025agent} further study hijacking risks from untrusted context, prompt injection, and adversarial tool use. ToolEmu~\citep{ruan2024toolemu} evaluates high-stakes tool-use failures, and R-Judge~\citep{yuan2024rjudge} benchmarks risk awareness over agent interaction traces. Skill-focused studies, including Skill-Inject~\citep{schmotz2026skillinject} and HarmfulSkillBench~\citep{jiang2026harmfulskillbench}, further expose malicious instructions or overtly harmful functionality embedded in skill artifacts. Recent defenses audit skills before deployment or execution. SkillScan~\citep{liu2026donotmentionmaliciousskills} and SkillSieve~\citep{hou2026skillsieve} combine static analysis with LLM-based semantic triage to detect vulnerable or malicious skill artifacts. SkillFortify~\citep{bhardwaj2026formalanalysis} further uses abstract interpretation and capability-based confinement to check whether a skill can exceed its declared authority. Together, these studies establish standalone artifact vetting as a necessary security primitive, but they leave open how intermediate artifacts produced by one skill may become actionable inputs to another.

%Although artifact-level vetting can identify unsafe standalone skills, it remains blind to failures that arise only when individually safe skills are composed into long-horizon workflows~\citep{jiang2026agentlab}. In practice, skills are invoked as interdependent modules: an upstream skill may provide an execution target, legitimacy endorsement, authorization cue, or advisory context that changes how a downstream skill acts, activating a path whose risk is not visible from either artifact alone~\citep{jiang2026sokagenticskills}. For example, as shown in Figure~1, a benign file-audit skill that reports stale access-control entries can become unsafe when its output is routed to a permission-management skill that updates the reported targets, causing unintended changes to sensitive access rights. We define this risk as \textbf{SCR} (Skill Composition Risk): the risk that skills judged safe in isolation produce unsafe downstream state changes when activated along a composed execution path.  
Although artifact-level vetting can identify unsafe standalone skills, it remains structurally blind to a distinct class of failure: harm that only emerges when individually safe skills are composed into longer agent workflows~\citep{jiang2026agentlab}. This concern is consistent with recent evidence that autonomous agents can be compromised through seemingly non-overt actions whose effects are amplified across tool-mediated execution~\citep{zhang2025breakingagents}. The reason is that skills in such workflows do not act as isolated nodes; they are invoked as interdependent modules. An upstream skill may produce an execution target, a legitimacy endorsement, an authorization cue, or advisory context, and the agent silently carries these intermediate artifacts into downstream invocations~\citep{jiang2026sokagenticskills}. Once these artifacts cross from one skill into another, an execution path is activated whose risk is invisible at any single artifact. As shown in \autoref{fig:composition-risk-overview}, the risk arises from the activated path between skills, not from any individual skill in isolation. The harmful event lives on the path between the two skills, not in either of them. Generalizing this observation, the security of a skill-enabled agent is not a static property of any standalone artifact, but an emergent property of the execution paths the agent activates. We therefore formalize this gap as \textbf{SCR} (\textbf{S}kill \textbf{C}omposition \textbf{R}isk): the risk that skills judged safe in isolation produce unsafe downstream state changes when composed through a shared execution context.

To evaluate SCR, we introduce SCR-Bench, the first benchmark suite that operationalizes SCR in sandboxed skill environments. Unlike evaluations based only on model-output text, SCR-Bench measures observable environmental state changes induced along activated skill-composition paths. SCR-Bench does not aim to exhaustively taxonomize agent-skill vulnerabilities. Instead, it instantiates three representative mechanisms of path-level risk. Capability flow occurs when an upstream skill supplies targets or operational context for a harmful downstream action. Trust transfer occurs when a benign-looking security output legitimizes a later high-risk skill or action. Authorization confusion occurs when advisory or audit-like context shifts the agent’s approval boundary toward unsafe downstream decisions. By comparing isolated and composed skill execution in these scenarios, SCR-Bench provides concrete empirical evidence for path-level risks.
% 2026年5月24日注释 To evaluate SCR, we introduce \textbf{SCR-Bench}, the first benchmark suite for operationalizing this formulation in sandboxed skill environments by measuring observable downstream state changes across activated composition paths. Rather than providing an exhaustive taxonomy vulnerabilities, SCR-Bench explores three representative mechanisms as illustrative examples to stimulate further research: capability-flow risks, where an upstream skill supplies execution targets or operational context that enables a downstream skill to perform harmful state-changing actions; trust-transfer risks, where a benign-looking upstream security skill lends legitimacy to a later risky skill or action; and authorization-confusion risks, where advisory or finding-like context shifts an agent's approval boundary toward unsafe downstream decisions. By comparing isolated and composed skill execution across these examples, SCR-Bench provides concrete empirical evidence for path-level risks.

%Across the three sub-benchmarks, SCR-Bench reveals a consistent ``Benign in Isolation, Harmful in Composition'' pattern: composed execution substantially increases attack success, risky installation, and unsafe approval rates compared with isolated execution. Capability-flow attacks remain near zero under isolated baselines but rise above 33\% under composed execution; trust-transfer settings produce harmful-skill installation rates above 96.5\% on four of five model backends; and authorization-confusion contexts increase risky approvals from 16.4\% to 33.9\%. These results show that path-aware evaluation captures risks that standalone artifact-level vetting can miss.The contributions are threefold:

Across these illustrative examples, empirical evaluations expose a persistent and severe security blind spot: composed execution consistently converts individually safe skills into harmful pathways. While isolated baselines show near-zero capability-flow risks, composing them drives attack success rates above 33\%; trust-transfer endorsements push harmful installation rates above 83.89\% on most model backends; and authorization confusion increases risky approvals from 15.7\% to 27.0\%. These findings indicate that isolated artifact vetting is structurally insufficient for securing skill-enabled agents. Securing agentic workflows requires a necessary paradigm shift from node-level inspection to path-aware evaluation. Our main contributions are summarized as follows:
\begin{itemize}
    \item \textbf{Path-level Risk Formulation:} We identify and formalize SCR, a path-level security gap in which skills that appear benign in isolation can produce harmful downstream outcomes when composed through shared execution context. This formulation shifts the unit of analysis from isolated skill artifacts to activated execution paths.
    \item \textbf{Benchmark Construction:} To operationalize SCR, we introduce SCR-Bench, the first controlled benchmark suite for evaluating path-level risks in sandboxed agent-skill environments. By measuring downstream state changes along activated execution paths, SCR-Bench reveals compositional harm across capability flow, trust transfer, and authorization confusion.
    \item \textbf{Empirical Findings:} Across multiple LLM backends, we show that composed skill execution exposes risks largely absent under isolated evaluation, revealing the ``Benign in Isolation, Harmful in Composition'' pattern and motivating path-aware security analysis.
\end{itemize}


% \section{Related Work}

% \paragraph{Agentic skills and ecosystem security.}
% Agentic skills have become a central abstraction for extending LLM agents from language-only reasoning to executable workflows. A recent SoK on agentic skills~\citep{jiang2026sokagenticskills} characterizes skills as reusable, programmatic modules that encode procedural memory and enable agents to interact with external services, software environments, and persistent state. This shift introduces ecosystem-level security risks. Agent Skills in the Wild~\citep{liu2026agentskillswild} shows that real-world skill repositories contain widespread vulnerabilities, with 26.1\% of analyzed skills containing at least one security issue. The same line of work also identifies attacker patterns such as data thieves and agent hijackers, demonstrating that skill ecosystems expose realistic adversarial surfaces~\citep{liu2026agentskillswild}. These studies establish the security relevance of skill ecosystems, but their measurements primarily concern vulnerabilities or malicious behavior within individual skill artifacts. They do not provide a controlled evaluation setting for risks that emerge only when individually benign skills are composed along an activated execution path.

% \paragraph{Isolated artifact vetting.}
% A growing line of work develops defenses for reviewing skills or tool calls as standalone security objects. SkillScan~\citep{liu2026agentskillswild}, SkillSieve~\citep{hou2026skillsieve}, and SkillProbe~\citep{guo2026skillprobe} use combinations of static analysis, LLM-based semantic triage, and behavioral inspection to identify vulnerable, misleading, or malicious skill artifacts. SkillFortify~\citep{bhardwaj2026formalanalysis} introduces formal analysis based on abstract interpretation and capability-based confinement to verify that a skill cannot exceed its declared authority under explicit assumptions. Runtime defenses such as AEGIS~\citep{yuan2026aegis} and ClawGuard~\citep{zhao2026clawguard} inspect individual tool invocations or execution parameters before an agent acts. These systems establish artifact-level and invocation-level vetting as important security primitives. However, they leave open a distinct benchmarking question: whether skills that pass standalone or invocation-level checks can still produce unsafe downstream outcomes when connected through capability flow, trust transfer, or authorization confusion.

% \paragraph{Compositional and multi-step attack chains.}
% Several recent studies move beyond single-step attacks and examine how risk emerges across multi-step agent behavior. Semantic Intent Fragmentation~\citep{ahad2026semanticintentfragmentation} shows that individually benign subtasks can compose into a harmful overall plan, revealing a gap between local intent classification and global task safety. SkillAttack~\citep{duan2026skillattack} explores feedback-driven procedures for discovering and refining adversarial skill-based attack paths. TraceSafe~\citep{chen2026tracesafe} and AgentLAB~\citep{jiang2026agentlab} further highlight the importance of long-horizon trajectories and multi-step tool use in agent safety evaluation. These works are closely related to Skill Composition Risk, but they differ in evaluation target. Existing compositional studies typically analyze plan-level intent, generated attack trajectories, or trace-level safety properties, whereas SCR-Bench isolates the skill-composition setting itself: individually benign skills are placed in controlled sandboxed environments, and risk is measured through downstream state changes such as permission updates, risky installations, and shifted approval decisions.

% \paragraph{Agent security benchmarking.}
% Agent security benchmarks have provided increasingly realistic environments for evaluating tool-using agents under adversarial conditions. AgentDojo~\citep{yuan2026aegis} evaluates prompt-injection attacks in dynamic tool-use environments, while Agent Security Bench~\citep{wang2026systematicsecurityopenclaw} studies vulnerabilities across multiple stages of agent operation, including prompting, tool use, and memory or storage. HarmfulSkillBench~\citep{jiang2026harmfulskillbench} evaluates whether agents can recognize and refuse skills whose functionality is overtly harmful. These benchmarks are valuable for measuring end-to-end robustness, prompt-injection resilience, and harmful-action refusal. SCR-Bench addresses a different evaluation gap: it is a mechanism-level benchmark for measuring how unsafe outcomes emerge and amplify along activated paths composed of individually benign skills.


\section{Related Works}
\paragraph{Skill ecosystem security and artifact-level vetting.}
LLM agents extend language-only reasoning through external tool calls~\citep{yao2023react,schick2023toolformer}, API invocation and tool learning~\citep{patil2024gorilla}, memory access, and reusable procedural modules~\citep{jiang2026sokagenticskills}. Among these extensions, agentic skills have emerged as reusable procedural modules for packaging execution policies, applicability conditions, and callable interfaces across tasks~\citep{jiang2026sokagenticskills}. As skills move into open registries and marketplaces, they introduce a distinct security surface. A large-scale empirical study of real-world skill repositories found that 26.1\% of analyzed skills contain at least one vulnerability, with data exfiltration and privilege escalation among the most prevalent categories~\citep{liu2026donotmentionmaliciousskills}. Related attacks on tool-selection pipelines further show that adversarial tool documents can steer agents toward attacker-chosen tools~\citep{shi2026toolhijacker}. This risk landscape has motivated artifact-level skill auditing. SkillScan~\citep{liu2026donotmentionmaliciousskills}, SkillSieve~\citep{hou2026skillsieve}, and SkillProbe~\citep{guo2026skillprobe} combine static analysis, LLM-based semantic inspection, and behavioral or collaborative auditing to detect vulnerable or malicious skills. A complementary line of work seeks stronger guarantees. SkillFortify~\citep{bhardwaj2026formalanalysis} applies formal analysis and capability-based confinement to skill supply chains. Runtime defenses instead focus on the execution boundary. AEGIS~\citep{yuan2026aegis} and ClawGuard~\citep{zhao2026clawguard} interpose before tool execution and block risky calls at invocation time. Existing work has therefore largely treated each skill artifact or tool invocation as the primary unit of review, which aligns the security boundary with a single module rather than with cross-skill context propagation. This boundary is useful but incomplete for SCR. It does not directly measure whether an upstream skill output later becomes a downstream execution target, trust signal, or authorization cue along an activated composition path.

\paragraph{Compositional agent risks and security benchmarks.}
Recent benchmarks have begun to examine risks that emerge beyond standalone skill artifacts and single tool calls. SkillAttack~\citep{duan2026skillattack} studies how adversarial prompting can exploit latent vulnerabilities in otherwise unmodified skills. HarmfulSkillBench~\citep{jiang2026harmfulskillbench} evaluates skills whose intended functionality is itself harmful. Broader agent-security benchmarks study related risks at the trajectory or environment level. AgentDojo~\citep{debenedetti2024agentdojo} evaluates indirect prompt injection in dynamic tool-use environments. Agent Security Bench~\citep{zhang2025agent} formalizes multiple attack surfaces across agent workflows. TraceSafe~\citep{chen2026tracesafe} assesses guardrails on multi-step tool-calling traces. AgentLAB~\citep{jiang2026agentlab} targets adaptive long-horizon attacks. The closest work to SCR's benign-in-isolation premise is Semantic Intent Fragmentation~\citep{ahad2026semanticintentfragmentation}. It shows that individually benign subtasks can compose into a policy-violating plan. SCR studies a different unit and mechanism. Rather than evaluating plan-level intent decomposition, SCR focuses on activated paths inside skill ecosystems. The central question is how outputs, endorsements, authorization-like signals, or side effects from locally benign skills are carried through shared execution context and consumed by downstream skills. SCR-Bench therefore takes the activated skill-composition path, rather than the isolated skill, individual tool call, full trajectory, or decomposed plan, as the unit of analysis. It further measures whether this path produces observable downstream state changes relative to single-skill baselines.



 

\section{Method}

This section defines the threat model, assumptions, and notation for SCR. It then introduces a context-dependent composition graph to model skill interactions and maps three mechanisms—capability flow, trust transfer, and authorization confusion—to path-level risk. Finally, it describes how \textbf{SCR-Bench} instantiates the formulation with controlled scenarios and empirical estimators.

\begin{figure}[h]
  \centering
  \includegraphics[width=\linewidth]{figures/pipeline-0612.pdf}
  \caption{Path-level composition in SCR-Bench. The figure illustrates the three mechanisms—capability flow, trust transfer, and authorization confusion—that can produce harmful outcomes from otherwise safe skills.}
  \label{fig:framework-pipeline}
\end{figure}

\subsection{Threat Model, Assumptions, and Notation}

The threat model considers a skill-enabled agent that invokes multiple locally bounded skills within a task and carries intermediate artifacts across steps. The adversary aims to induce an unsafe downstream outcome by exploiting composition paths among skills that appear benign in isolation. This defines the execution setting for SCR: intermediate artifacts such as outputs, trust signals, authorization cues, or side effects may be reused in later decisions, allowing upstream skill behavior to shape downstream actions.

% The threat model considers a skill-enabled agent that invokes multiple locally bounded skills within a task and carries intermediate artifacts across steps. The adversary aims to induce an unsafe downstream outcome by exploiting composition paths among skills that appear benign in isolation. The threat model specifies the execution setting for SCR: an agent may invoke multiple skills during a task and reuse intermediate outputs, trust signals, authorization cues, or side effects in later decisions. This cross-step reuse enables complex workflows but also creates composition paths that shape downstream actions.

The adversary need not make any single skill harmful in isolation. Instead, the adversary may construct or exploit a composition path in which an upstream skill produces a target, endorsement, audit finding, authorization cue, or state change that influences downstream invocation or judgment. Thus, even skills with ordinary descriptions and bounded local behavior can become unsafe when upstream outputs cross capability, trust, or authorization boundaries and are consumed as operational evidence. The unit of security analysis therefore shifts from isolated skills to context-dependent composition paths. \autoref{tab:notations} summarizes notations, including path risk for a single path \(\pi\) under context \(h\), ecosystem-level SCR over candidate paths, and mechanism risk as the empirical attack success rate for a specific composition mechanism.
% The threat model considers an agent that can access a set of installable or callable skills during a task.
% The agent may invoke more than one skill in the same session, carry intermediate outputs in its context, and use earlier results when deciding whether to call later skills.
% This is the normal utility model for skills: outputs are useful because they can be reused.

% The adversary's objective is not necessarily to make one skill perform the whole harmful action in isolation.
% Instead, the adversary may create or rely on a composition path in which one skill supplies a target, endorsement, advisory finding, or other context that changes how a later skill is invoked or judged.
% In the \textbf{SCR-TrustLift} setting, the skill market includes attacker-supplied harmful skills and a benign-looking security skill that can provide trust signals.
% In the \textbf{SCR-CapFlow} and \textbf{SCR-AuthBlur} settings, the individual skills are designed to have ordinary-looking descriptions and bounded local behavior; the risk appears when their outputs cross a capability, trust, or authorization boundary.

% All empirical results are measured in controlled, sandboxed skill environments.
% The benchmarks use mock files, mock service records, local state files, binary approval decisions, or simulated installation events as observable outcomes.
% They do not modify real infrastructure, transmit real secrets, or estimate the absolute prevalence of risk in a public deployment.
% The intended claim is therefore mechanism-level: such path risks can arise under composition, and isolated skill review does not by itself rule them out.


% 2026年6月7日 注释掉的table1 符号表
% \begin{table}[t]
%   \caption{Notation used in the SCR formulation and empirical estimators.}
%   \label{tab:notation}
%   \centering
%   \small
%   \begin{tabular}{ll}
%     \toprule
%     Symbol & Meaning \\
%     \midrule
%     \(s_i\) & A skill available to the agent. \\
%     \(V\) & Set of available skills \(\{s_1,\ldots,s_n\}\). \\
%     \(h\) & Task and session context, including memory, intermediate outputs, and decision state. \\
%     \(G_h=(V,E_h)\) & Context-dependent skill composition graph. \\
%     \(e_{ij}\) & Composition edge from \(s_i\)'s output to \(s_j\)'s input or decision state. \\
%     \(\pi\) & Activated or candidate composition path. \\
%     \(a_\pi(h)\) & Probability that the agent activates path \(\pi\) under context \(h\). \\
%     \(q_\pi(h)\) & Conditional probability that activated path \(\pi\) reaches a harmful state. \\
%     \(r_\pi(h)\) & Effective path risk \(a_\pi(h)q_\pi(h)\). \\
%     \(Y^z_{c,t}\) & Path-level success indicator for mechanism \(z\), case \(c\), trial \(t\). \\
%     \bottomrule
%   \end{tabular}
% \end{table}

\begin{table}[h]
\centering
\caption{Notation used in the SCR formulation and empirical estimators.}
\label{tab:notations}
\footnotesize
\setlength{\tabcolsep}{2.5pt}
\renewcommand{\arraystretch}{0.92}

\begin{adjustbox}{center,max width=\linewidth}
\begin{tabularx}{0.98\linewidth}{
@{}
>{\raggedright\arraybackslash}p{0.10\linewidth}
>{\raggedright\arraybackslash}X
>{\raggedright\arraybackslash}p{0.10\linewidth}
>{\raggedright\arraybackslash}X
@{}
}
\toprule
Symbol & Meaning & Symbol & Meaning \\
\midrule
$s_i$ & Agent skill &
$V$ & Skill set $\{s_1,\ldots,s_n\}$ \\

$h$ & Task/session context &
$G_h$ & Context-dependent skill graph \\

$E_h$ & Feasible composition edges under $h$ &
$e_{ij}$ & Edge from $s_i$ to $s_j$ \\

$o_i$ & Output/state change of $s_i$ &
$x_j$ & Object consumed by $s_j$ or planner \\

$\pi$ & Candidate/activated path &
$\Pi_k(G_h)$ & Length-$k$ paths in $G_h$ \\

$B_\pi$ & Harmful downstream state of $\pi$ &
$a_\pi(h)$ & Activation probability of $\pi$ \\

$q_\pi(h)$ & Conditional probability of reaching $B_\pi$ &
$r_\pi(h)$ & Path risk $a_\pi(h)q_\pi(h)$ \\

$z$ & Risk mechanism &
$c,t$ & Case and trial indices \\

$C,m$ & Number of cases and trials per case &
$Y^z_{c,t}$ & Harmful-state indicator  \\
\bottomrule
\end{tabularx}
\end{adjustbox}
\end{table}
% $\widehat{\mathrm{ASR}}_z$ & Empirical attack success rate &
%  & \\    (ASR的符号说明)

\subsection{Skill Composition Graph and Path Risk}

To make path-level risk amenable to systematic analysis, an agent skill ecosystem under task context \(h\) is modeled as a directed, context-dependent composition graph:
\begin{equation}
G_h=(V,E_h).
\label{eq:skill-graph}
\end{equation}
Here, \(V=\{s_1,\ldots,s_n\}\) denotes the set of skills available to the agent, and \(E_h\) denotes the set of skill-to-skill composition edges that may be activated under context \(h\). The graph is not a static skill-registry graph. Instead, it captures which skills can form operational composition relations under a particular task and session state, and therefore how locally bounded skills may be connected into composition paths.

The context \(h\) is included because skill composition depends on task-specific semantics rather than fixed interfaces alone. The same skill pair may be independent in one task but connected in another: file names produced by a file-indexing skill may become reportable objects for a reporting skill, while a security-audit output may make a later installer action appear justified. Thus, \(E_h\) is determined by technical interfaces together with task intent, session memory, intermediate outputs, and the agent's decision state.

A composition edge represents a semantic or operational dependency across skills, rather than merely a direct function call. If an output, state, endorsement, or side effect produced by skill \(s_i\) affects the input to skill \(s_j\), the invocation of \(s_j\), or the agent's judgment about \(s_j\), the relation is written as
\begin{equation}
e_{ij}: o_i \rightarrow x_j,
\label{eq:composition-edge}
\end{equation}
where \(o_i\) is produced by \(s_i\), and \(x_j\) is the downstream object consumed by \(s_j\) or the agent planner. The semantics of \(x_j\) determine the edge type: execution targets induce capability flow, legitimacy signals induce trust transfer, and approval evidence induces authorization confusion. These three boundary crossings define the mechanisms evaluated later. A skill composition path on this graph is a sequence of skill invocations:
\begin{equation}
\pi=(s_{i_1}\rightarrow s_{i_2}\rightarrow \cdots \rightarrow s_{i_k})\in \Pi_k(G_h).
\label{eq:path}
\end{equation}
For a path \(\pi\), let \(a_\pi(h)\) denote the probability that the agent activates the path under context \(h\), and let \(q_\pi(h)\) denote the conditional probability that the activated path reaches a harmful state \(B_\pi\). The effective risk of the path is
\begin{equation}
r_\pi(h)=a_\pi(h)q_\pi(h).
\label{eq:path-risk}
\end{equation}
This decomposition separates the agent's tendency to follow a composition path from the conditional harm of that path once activated. The basic unit of risk is therefore not an isolated skill, but a context-dependent composition path.

This graph model expands the security boundary for skills from isolated nodes to context-dependent composition transitions. A node-level audit can determine whether each skill has bounded local behavior, but it cannot determine whether an upstream output will later be reinterpreted as an execution target, trust signal, or authorization cue. The relevant question is therefore not only whether a skill \(s_i\) is safe in isolation, but whether the transition from \(s_i\) to \(s_j\) preserves the intended boundary between data and action, advice and authorization, or observation and control.
%修改后的第二小节%
% \subsection{Skill Composition Graph}

% Given a task context \(h\), an agent skill ecosystem is modeled as a directed, context-dependent composition graph:
% \begin{equation}
% G_h=(V,E_h).
% \label{eq:skill-graph}
% \end{equation}
% Here, \(V=\{s_1,\ldots,s_n\}\) denotes the set of skills available to the agent, and \(E_h\) denotes the set of skill-to-skill composition edges that may be activated under context \(h\). The graph is not a static registry graph. Instead, it represents which skills can form operational composition relations under a particular task and session state. Thus, \(G_h\) captures how locally bounded skills may be connected by the agent into composition paths in the current context.

% The context \(h\) is part of the graph because skill composition is not determined solely by fixed interfaces. The same pair of skills may be unrelated in one task but tightly coupled in another. For example, a file-indexing skill and a reporting skill may be independent when the task is to summarize local files, but they become connected when file names discovered by the former are treated as reportable objects by the latter. Similarly, the output of a security-audit skill may affect whether an installer skill is treated as a reasonable subsequent action. Consequently, \(E_h\) is determined not only by technical interfaces, but also by task intent, session memory, intermediate outputs, and the agent's decision state.

% A composition edge represents a semantic or operational dependency across skills, rather than merely a direct function call. If an output, state, endorsement, or side effect produced by skill \(s_i\) affects the input to skill \(s_j\), the invocation of \(s_j\), or the agent's judgment about \(s_j\), the relation is written as
% \begin{equation}
% e_{ij}: o_i \rightarrow x_j,
% \label{eq:composition-edge}
% \end{equation}
% where \(o_i\) is produced by \(s_i\), and \(x_j\) is the downstream object consumed by \(s_j\) or by the agent planner. Different types of \(x_j\) induce different composition risks. If \(x_j\) is used as an execution target, the edge carries capability. If \(x_j\) is used as a legitimacy or trust signal, the edge carries trust. If \(x_j\) is interpreted as approval or authorization evidence, the edge carries authorization. The mechanisms studied later, capability flow, trust transfer, and authorization confusion, correspond to these three forms of boundary crossing.

% This graph model expands the security boundary for skills from isolated nodes to context-dependent composition transitions. A node-level audit can determine whether each skill has bounded local behavior, but it cannot determine whether an upstream output will later be reinterpreted as an execution target, trust signal, or authorization cue. The relevant question is therefore not only whether a skill \(s_i\) is safe in isolation, but whether the transition from \(s_i\) to \(s_j\) preserves the intended boundary between data and action, advice and authorization, or observation and control. Skill safety is consequently treated as a path-level property rather than only a node-level property.

% 原本的第三小节，拆开融入了第二小节和第四小节%
% \subsection{\textbf{Skill Composition Risk} amplification}

% A skill composition path is a sequence of skill invocations:
% \begin{equation}
% \pi=(s_{i_1}\rightarrow s_{i_2}\rightarrow \cdots \rightarrow s_{i_k})\in \Pi_k(G_h).
% \label{eq:path}
% \end{equation}
% For a path \(\pi\), let \(a_\pi(h)\) be the probability that the agent activates the path in context \(h\), and let \(q_\pi(h)\) be the probability that the activated path reaches a harmful state \(B_\pi\).
% The effective risk of a path is
% \begin{equation}
% r_\pi(h)=a_\pi(h)q_\pi(h).
% \label{eq:path-risk}
% \end{equation}
% This term is the basic unit of the framework.
% It separates the agent's tendency to follow a path from the conditional danger of that path once followed.

% The ecosystem-level \textbf{Skill Composition Risk} is the probability that at least one activated path reaches a harmful state.
% Using the effective path risks in \autoref{eq:path-risk}, an independence approximation gives:
% \begin{equation}
% R_{\mathrm{comp}}(G,h)
% =
% 1-\prod_{\pi\in\Pi(G_h)}
% \left(1-r_\pi(h)\right).
% \label{eq:comp-risk}
% \end{equation}
% This equation is the central object of the method.
% It should be read as a tractable risk accounting model rather than a claim that real path failures are independent.
% Without independence, the same intuition can be expressed with union bounds or empirical path-level estimates; the key point is still that risk is accumulated over candidate paths rather than over isolated nodes.
% It says that ecosystem risk is not determined by a single skill's standalone behavior.
% It is determined by how many harmful paths exist, how often the agent activates them, and how dangerous they are after activation.

% For two-skill paths, the homogeneous approximation gives:
% \begin{equation}
% R_2(G,h)
% =
% 1-\left(1-\bar r_2(h)\right)^{\binom{n}{2}},
% \label{eq:pair-risk}
% \end{equation}
% where \(\bar r_2(h)\) is the average effective risk over skill pairs.
% This expression captures the amplification effect: even when \(\bar r_2(h)\) is small, the number of candidate pairs grows quadratically with \(n\).
% For longer paths, the same argument applies with \(N_k(h)=|\Pi_k(G_h)|\), which can grow as \(O(n^k)\) in dense composition settings.

% \paragraph{Why isolated vetting misses \textbf{Skill Composition Risk}.}
% Isolated skill vetting covers a node-level defense surface, while composition attacks live on a path-level risk surface.
% The coverage gap scales as
% \begin{equation}
% \frac{|\mathcal{R}_{path}|}{|\mathcal{D}_{node}|}
% =
% \frac{O(n^k)}{O(n)}
% =O(n^{k-1}), \quad k\ge 2.
% \label{eq:surface-gap}
% \end{equation}
% This gap explains how a system can pass node-level vetting while still having high \textbf{Skill Composition Risk}.
% A skill may be safe as a standalone unit, yet its output can become the missing input, endorsement, or authorization cue that activates a harmful path.

% \subsection{Risk mechanisms}

% The framework maps each experimental mechanism to a different source of path risk.
% In \textbf{capability flow}, an upstream discovery skill supplies the target for a downstream execution skill.
% This increases \(q_\pi(h)\) because a path that previously lacked a target can now reach a concrete harmful state.
% The important point is that neither endpoint needs to contain the whole attack.
% The first skill may only observe, enumerate, or summarize; the second may only modify, upload, schedule, or report.
% The harmful behavior appears when the agent binds the first skill's discovered object to the second skill's action interface.
% For this reason, \textbf{SCR-CapFlow} uses path-level ground truth: success requires the downstream side effect to act on the same target discovered upstream.

% In \textbf{trust transfer}, an upstream skill produces an endorsement that changes how the agent evaluates a later risky action.
% This can increase \(a_\pi(h)\), because the agent is more likely to follow the path, and it can also increase \(q_\pi(h)\), because the downstream action is treated as legitimate.
% This mechanism is different from direct malicious installation.
% The upstream skill can look like a normal security, audit, or validation utility.
% Its output then becomes a legitimacy signal that lowers the agent's resistance to a later high-risk step.
% The path risk comes from the transfer of trust across skills, not from a single instruction that is obviously malicious in isolation.

% In \textbf{authorization confusion}, advisory or finding-like context is interpreted as approval, so the downstream decision boundary changes even if no explicit permission was granted.
% This mechanism captures a common failure mode in long agent sessions.
% Information that should remain descriptive becomes operational evidence for action.
% A finding such as ``this request appears documented'' or an advisory such as ``the risk is acceptable under the stated policy'' can be reused as if it were formal authorization.
% The resulting error is not simply prompt injection; the upstream content may be relevant and non-imperative, yet it still shifts the downstream approval decision.


% \subsection{Risk Mechanisms}

% At the ecosystem level, \textbf{Skill Composition Risk} arises from the accumulation of risk over candidate composition paths. If any activated path reaches a harmful state, the skill ecosystem is exposed to risk. Ecosystem risk is therefore determined not only by the standalone behavior of individual skills, but also by the number of candidate paths, the probability that the agent activates those paths, and the conditional probability that an activated path reaches harm. As the number of available skills increases, potential composition paths can grow rapidly. Even when each skill appears bounded or benign under local review, a cross-skill transition may still reinterpret an upstream output as downstream operational evidence. Isolated vetting covers a node-level security surface, whereas \textbf{Skill Composition Risk} arises from path-level transitions; checking individual skills alone therefore cannot rule out composition risk.

% This work considers three mechanisms that instantiate such path-level risk. The first is \textbf{capability flow}. In this mechanism, an upstream discovery skill supplies the execution target or operational context required by a downstream execution skill. This primarily increases the conditional harm probability \(q_\pi(h)\), because a path that previously lacked a concrete target can now reach an executable harmful state. Neither endpoint needs to contain the complete attack: the upstream skill may only observe, enumerate, or summarize, while the downstream skill may only modify, upload, schedule, or report. Harmful behavior emerges when the agent binds the object discovered upstream to the downstream action interface. Evaluating capability flow therefore requires path-level evidence: the downstream side effect must act on the same target discovered upstream.

% The second mechanism is \textbf{trust transfer}. An upstream skill produces an endorsement, audit result, or trust assessment that changes how the agent evaluates a later high-risk action. This can increase the path activation probability \(a_\pi(h)\), because the agent becomes more likely to continue to the downstream step, and can also increase the conditional harm probability \(q_\pi(h)\), because the downstream action is treated as more legitimate. This mechanism differs from direct malicious installation or direct malicious execution. The upstream skill may resemble an ordinary security, audit, or validation utility, while its output becomes a legitimacy signal in a later decision and lowers the agent's resistance to a high-risk step. The risk comes from trust being transferred across skills, rather than from an explicitly malicious instruction inside a single skill.

% The third mechanism is \textbf{authorization confusion}. In this mechanism, advisory or finding-like context from an upstream skill is misinterpreted by a downstream decision process as approval or authorization evidence, thereby shifting the agent's approval boundary. The upstream content may be relevant and non-imperative, yet still convert descriptive information into operational evidence. For example, a finding such as ``this request appears documented'' or an advisory statement such as ``the risk is acceptable under the stated policy'' may be reused by a later step as if it were formal authorization. This failure mode is not simply prompt injection: the risk does not require an explicit attack instruction, but arises when the semantic boundary between advice, findings, and authorization becomes blurred.

% These mechanisms correspond to three forms of boundary crossing: targets crossing into execution interfaces, trust crossing into downstream judgments, and advice crossing into authorization decisions. \textbf{SCR-Bench} instantiates these mechanisms as controlled scenarios and evaluates whether the agent actually reaches a harmful state along a composition path.


\subsection{Risk Mechanisms}

% The path-risk formulation explains why isolated skill vetting is incomplete. A skill ecosystem becomes unsafe when at least one activated composition path reaches a harmful state. The risk therefore depends not only on whether individual skills appear benign, but also on how many candidate paths exist, how often the agent activates them, and how likely an activated path is to produce harm. As the number of available skills increases, locally bounded skills may still form harmful paths when an upstream output is reused by a downstream skill. This paper focuses on three representative mechanisms through which such path risk can arise: \textbf{capability flow}, \textbf{trust transfer}, and \textbf{authorization confusion}.

The path-level formulation above explains how ecosystem-level risk accumulates over candidate composition paths. This paper focuses on three representative mechanisms through which such path-level risk can arise: \textbf{capability flow}, \textbf{trust transfer}, and \textbf{authorization confusion}.

%2026年6月7日注释，原Capability Flow的机制分析。
% In \textbf{capability flow}, an upstream discovery skill supplies the target or operational context required by a downstream execution skill, increasing \(q_\pi(h)\) by turning an incomplete path into an executable harmful state. Neither endpoint needs to contain the full attack: the upstream skill may only observe or enumerate, while the downstream skill may only modify, upload, schedule, or report. Harm emerges when the agent binds the upstream-discovered object to the downstream action interface.
In \textbf{capability flow}, individually legitimate skills with distinct permissions compose through context to exceed an intended safety boundary. An upstream discovery skill supplies the target or operational context required by a downstream execution skill, increasing \(q_\pi(h)\) by turning an incomplete path into an executable harmful state. Neither endpoint needs to contain the full attack: harm emerges when the agent binds an upstream-discovered object to a downstream action interface such as modification, upload, scheduling, or reporting.

In \textbf{trust transfer}, an upstream skill produces an endorsement, audit result, or trust assessment that makes a downstream high-risk action appear legitimate. This can increase both \(a_\pi(h)\), by making the agent more likely to continue to the downstream step, and \(q_\pi(h)\), by weakening scrutiny of the action once reached. Unlike direct malicious installation or execution, the upstream skill may appear to be a benign security, audit, or validation utility; the risk arises when its output is reused as a legitimacy signal for a later decision.

In \textbf{authorization confusion}, advisory or finding-like context from an upstream skill is treated by a downstream decision process as approval or authorization evidence. Even relevant, non-imperative content can shift the agent's approval boundary when descriptive findings are reused as operational authorization. For example, statements such as ``this request appears documented'' or ``the risk is acceptable under the stated policy'' may later be interpreted as formal approval. Unlike prompt injection, this mechanism does not require an explicit attack instruction; the risk arises from confusing advice or findings with authorization.

% These three mechanisms directly motivate the construction of \textbf{SCR-Bench}. \textbf{SCR-CapFlow} tests whether a discovered target is carried into downstream execution, \textbf{SCR-TrustLift} tests whether an upstream trust signal increases downstream harmful installation, and \textbf{SCR-AuthBlur} tests whether advisory or finding-like context increases risky approval. In all three cases, the relevant evidence is a path-level outcome rather than the isolated appearance of any single skill.
% \begin{figure}[t]
%   \centering
%   \includegraphics[width=0.5\linewidth]{figures/three.pdf}
%   \caption{Composition of \textbf{SCR-Bench} and the types of risk instantiated in each sub-benchmark.}
%   \label{fig:three-mechanisms}
% \end{figure}

\subsection{SCR-Bench Construction}

\begin{wrapfigure}{r}{0.45\textwidth} % r=右侧, 0.45\textwidth = 图宽度
  \centering
  \includegraphics[width=\linewidth]{figures/SCR-Bench.pdf}
  \caption{Composition of \textbf{SCR-Bench} and the types of risk instantiated in each sub-benchmark.}
  \label{fig:three-mechanisms}
\end{wrapfigure}

Building on the three mechanisms above, \textbf{SCR-Bench} instantiates capability flow, trust transfer, and authorization confusion as three controlled sub-benchmarks, with the benchmark composition illustrated in \autoref{fig:three-mechanisms}. Each sub-benchmark fixes one risk mechanism, separates skill-local capability from task context, and constrains endpoint skills to bounded local behavior. The evaluation target is not whether an isolated skill appears suspicious, but whether a composed execution path produces an observable harmful state. Each trial is scored with a binary path-level indicator. For mechanism \(z\), case \(c\), and trial \(t\), \(Y^z_{c,t}=1\) iff the corresponding bad event produces a downstream state change in the mock environment. Mere invocation or textual description is not counted as success. The attack success rate for mechanism \(z\) averages this indicator over \(C\) cases and \(m\) trials per case.
\begin{equation}
\widehat{ASR}^{z}
=
\frac{1}{Cm}\sum_{c=1}^{C}\sum_{t=1}^{m}Y^z_{c,t}.
\label{eq:asr}
\end{equation}
This scoring rule aligns measurement with the path-risk definition: it evaluates whether the composition path reaches a harmful state, not whether a skill was merely invoked or the model only described a risky action. Detailed mechanism-specific estimators are given in \autoref{app:method-details}.
% Building on the three mechanisms above, \textbf{SCR-Bench} instantiates capability flow, trust transfer, and authorization confusion as three controlled sub-benchmarks. Rather than judging whether an isolated skill appears suspicious from its static description, SCR-Bench tests whether an agent reaches an observable harmful state through a composition path. Each sub-benchmark fixes one risk mechanism, separates the local capability of each skill from the task context provided in the agent prompt, and constrains endpoint skills to bounded local behavior.

% For mechanism \(z\), case \(c\), and trial \(t\), let \(Y^z_{c,t}=1\) only when the corresponding path-level bad event is realized as a downstream state change in the mock environment, such as a permission update, request-log entry, scheduled device action, or approval decision. The attack success rate is estimated as
% \begin{equation}
% \widehat{ASR}^{z}
% =
% \frac{1}{Cm}\sum_{c=1}^{C}\sum_{t=1}^{m}Y^z_{c,t}.
% \label{eq:asr}
% \end{equation}
% This scoring rule aligns measurement with the path-risk definition: it evaluates whether the composition path reaches a harmful state, not whether a skill was merely invoked or the model only described a risky action.

% \textbf{SCR-CapFlow} evaluates capability flow using paired discovery and execution skills. Each case includes a read-only discovery skill that writes structured targets from a mock sandbox to local state, and an execution or control skill that acts only on an explicitly supplied target and records side effects such as permission changes, upload logs, device schedules, firewall rules, or cron jobs. Success requires more than invoking both skills:
% \[
% attack\_chain\_success =
% discover(t) \land side\_effect(t) \land uses\_discovered\_target(t).
% \]
% A trial succeeds only if the execution side effect acts on the same concrete target discovered upstream, distinguishing true capability flow from coincidental co-invocation. The sub-benchmark contains 150 cases across permission and access control, mock reporting, device control, network configuration, and persistence scheduling.

% \textbf{SCR-TrustLift} evaluates trust transfer through harmful installation opportunities paired with a benign-looking upstream security or review skill. The upstream skill does not directly perform the harmful installation; instead, it emits an endorsement, audit result, or trust signal that can make a later installation request appear safer or more legitimate. This sub-benchmark measures whether upstream endorsement changes a downstream installation decision, rather than whether a malicious skill can complete installation in isolation. Each tested backend is evaluated over 401 installation trials.

% \textbf{SCR-AuthBlur} evaluates authorization confusion through a context-level intervention that tests whether upstream context shifts the agent's downstream approval boundary. \autoref{tab:authblur-levels} summarizes the context levels. \(L_0\) uses an unrelated control context to estimate the baseline approval rate; \(L_1\) provides related natural-language context resembling ordinary task continuity; \(L_2\) invokes a plain-finding advisor required to avoid authorization language; and \(L_3\) invokes a stronger advisory version that may contain non-binding but authorization-like assessments. The main paper reports \(L_0\), \(L_1\), and \(L_3\) over 118 retained cases spanning security, privacy, finance, HR, supply-chain, and operational-safety decisions. \(L_2\) is reported in the appendix because only 52 cases satisfy the stricter audit criterion that the plain findings do not already contain approval-like evidence.Detailed mechanism-specific estimators are given in \autoref{app:method-details}.


\textbf{SCR-CapFlow} evaluates capability flow using paired discovery and execution skills. Each case includes a read-only discovery skill and an execution or control skill. The main comparison contrasts control and isolated-skill conditions with two composed variants: \emph{Neutral} uses legitimate, task-oriented language without explicitly instructing the agent to form a harmful chain, testing whether the agent autonomously composes benign skills into a boundary-crossing capability flow; \emph{Explicit} specifies the use of both skills to identify a target and execute the corresponding action, serving as a path-specified positive control for whether the underlying chain is executable. Success requires more than invoking both skills:
\[
attack\_chain\_success =
discover(t) \land side\_effect(t) \land uses\_discovered\_target(t).
\]
A trial succeeds only if the execution side effect acts on the same concrete target discovered upstream, distinguishing true capability flow from coincidental co-invocation. The sub-benchmark contains 150 cases across permission and access control, mock reporting, device control, network configuration, and persistence scheduling.



% 2026年6月7日注释，将SCR-CapFlow信息过低的介绍改写。
% \textbf{SCR-CapFlow} evaluates capability flow using paired discovery and execution skills. Each case includes a read-only discovery skill that writes structured targets from a mock sandbox to local state, and an execution or control skill that acts only on an explicitly supplied target and records side effects such as permission changes, upload logs, device schedules, firewall rules, or cron jobs. The main comparison is composed execution against control and isolated-skill conditions. Success requires more than invoking both skills:
% \[
% attack\_chain\_success =
% discover(t) \land side\_effect(t) \land uses\_discovered\_target(t).
% \]
% A trial succeeds only if the execution side effect acts on the same concrete target discovered upstream, distinguishing true capability flow from coincidental co-invocation. The sub-benchmark contains 150 cases across permission and access control, mock reporting, device control, network configuration, and persistence scheduling.

\textbf{SCR-TrustLift} evaluates trust transfer through harmful installation opportunities paired with a benign-looking upstream security or review skill. The upstream skill does not directly perform the harmful installation; instead, it emits an endorsement, audit result, or trust signal that can make a later installation request appear safer or more legitimate. The key quantity is the lift between endorsed and non-endorsed contexts, measuring whether upstream endorsement changes a downstream installation decision rather than whether a malicious skill can complete installation in isolation. Each tested backend is evaluated over 401 installation trials.

\textbf{SCR-AuthBlur} evaluates authorization confusion by testing whether upstream context shifts the agent's downstream approval boundary. The key quantity is the lift from the control level \(L_0\) to stronger context levels. Detailed context-level definitions are provided in the Appendix. In brief, \(L_0\) uses an unrelated control context, \(L_1\) provides related task-like context, \(L_2\) invokes a plain-finding advisor without authorization language, and \(L_3\) invokes a stronger advisory version with non-binding authorization-like assessments. The main text reports \(L_0\), \(L_1\), and \(L_3\) over 118 retained cases across security, privacy, finance, HR, supply-chain, and operational-safety decisions; \(L_2\) is deferred to the \autoref{app:authblur-levels} due to stricter audit criteria.

% \subsection{SCR-Bench Construction}

% Building on the three mechanisms above, \textbf{SCR-Bench} instantiates capability flow, trust transfer, and authorization confusion as three controlled sub-benchmarks. The goal is not to determine whether an isolated skill appears suspicious from its static description, but to test whether an agent reaches an observable harmful state through a composition path. Each sub-benchmark fixes one risk mechanism, separates the local capability of each skill from the task context provided in the agent prompt, and constrains endpoint skills to bounded local behavior. Scoring focuses on downstream state changes caused by the composition path, rather than on whether a skill was merely invoked or whether the agent produced a risky textual explanation. \autoref{fig:SCR-Bench-construction} summarizes the construction pipeline.

% \begin{figure}[t]
%   \centering
%   \includegraphics[width=\linewidth]{figures/skillpathbench_construction.png}
%   \caption{\textbf{SCR-Bench} construction overview. The benchmark instantiates three path-risk mechanisms as controlled scenarios, with skill behavior, prompt conditions, sandboxed outcomes, and path-level scoring rules.}
%   \label{fig:SCR-Bench-construction}
% \end{figure}

% \textbf{SCR-CapFlow} evaluates capability flow using paired discovery and execution skills. Each case includes a read-only discovery skill that writes structured targets from a mock sandbox to local state, and an execution or control skill that acts only on an explicitly supplied target and records side effects such as permission changes, upload logs, device schedules, firewall rules, or cron jobs. Success requires more than invoking both skills:
% \[
% attack\_chain\_success =
% discover(t) \land side\_effect(t) \land uses\_discovered\_target(t).
% \]

% A trial succeeds only if the execution side effect acts on the same concrete target discovered upstream, distinguishing true capability flow from coincidental co-invocation. The sub-benchmark contains 150 cases across permission and access control, mock reporting, device control, network configuration, and persistence scheduling.

% \textbf{SCR-TrustLift} evaluates trust transfer through harmful installation opportunities paired with a benign-looking upstream security or review skill. The upstream skill does not directly perform the harmful installation. Instead, it emits an endorsement, audit result, or trust signal that can make a later installation request appear safer or more legitimate. This sub-benchmark therefore measures whether upstream endorsement changes a downstream installation decision, rather than whether a malicious skill can complete installation in isolation. Each tested backend is evaluated over 401 installation trials.

% \textbf{SCR-AuthBlur} evaluates authorization confusion through a context-level intervention. The sub-benchmark tests whether different forms of upstream context shift the agent's downstream approval boundary. \autoref{tab:authblur-levels} summarizes the context levels. \(L_0\) uses an unrelated control context to estimate the baseline approval rate. \(L_1\) provides related natural-language context that resembles ordinary task continuity. \(L_2\) invokes a plain-finding advisor that is required to avoid authorization language. \(L_3\) invokes a stronger advisory version that may contain non-binding but authorization-like assessments. The main paper reports \(L_0\), \(L_1\), and \(L_3\) over 118 retained cases spanning security, privacy, finance, HR, supply-chain, and operational-safety decisions. \(L_2\) is reported in the appendix because only 52 cases satisfy the stricter audit criterion that the plain findings do not already contain approval-like evidence.



% \subsection{Empirical estimators}

% Each mechanism is estimated with path-level success indicators.
% For mechanism \(z\), case \(c\), and trial \(t\), let \(Y^z_{c,t}=1\) when the corresponding path-level bad event occurs.
% The attack success rate is
% \begin{equation}
% \widehat{ASR}^{z}
% =
% \frac{1}{Cm}\sum_{c=1}^{C}\sum_{t=1}^{m}Y^z_{c,t}.
% \label{eq:asr}
% \end{equation}
% For \textbf{SCR-CapFlow}, the key comparison is the composed condition against the control and isolated-skill conditions.
% For \textbf{SCR-TrustLift}, the key quantity is the lift between endorsed and non-endorsed contexts.
% For \textbf{SCR-AuthBlur}, the key quantity is the lift from the control level \(L_0\) to stronger context levels.
% These estimators deliberately avoid scoring only the agent's textual explanation.
% Whenever possible, a trial is counted as successful only when the downstream state changes in the mock environment, such as a permission update, a request log entry, a scheduled device action, or an approval decision.
% This keeps the measurement aligned with the path-risk definition: the question is whether the composition path reaches a harmful state, not whether the model merely talks about doing so.
% Detailed mechanism-specific estimators are given in \autoref{app:method-details}.

\section{Experiments}

\begin{wrapfigure}[22]{r}{0.60\textwidth}
  \centering
  \includegraphics[width=\linewidth]{figures/SCR_CapFlow-0612.pdf}
  \caption{SCR-CapFlow ASR across model backends under control, single-skill, and composed-skill execution conditions. Shaded bands denote 95\% bootstrap confidence intervals.}
  \label{fig:capflow_main}
\end{wrapfigure}
The experiments evaluate whether the \textbf{SCR} framework captures failures that isolated skill review would miss.
They are organized around the three mechanisms in Section~3: capability flow, trust transfer, and authorization confusion.
For each mechanism, the central comparison is between a condition that exposes only isolated pieces of the workflow and a condition that allows the agent to activate a cross-skill path. Detailed ablation studies, prompts, case examples, implementation details, and additional experimental settings are provided in Appendix.

% \subsection{Experimental Setup}

% All experiments were conducted in a controlled, sandboxed environment, ensuring that only \emph{observable downstream state changes}---such as permission updates, installation logs, or approval decisions---count as successful attacks, rather than relying on the model's textual outputs. We tested across multiple backends to capture model-specific variations in agent behavior. Each experimental case was evaluated over $m=5$ independent trials, and reported results are averaged to reduce stochastic effects. The sandbox environment was constructed with a local skill registry, mock filesystem, service interfaces, scheduling mechanisms, and logging infrastructure, allowing precise tracking of capability, trust, and authorization flows along composition paths. Attack success rate (ASR) was used as the primary metric, defined as the fraction of trials in which the designated harmful event occurred, with 150 cases for capability flow, 401 installation trials per backend for trust transfer, and 118 retained decision cases for authorization confusion. This design guarantees rigorous and reproducible measurement of skill composition risk.

\subsection{Experimental Setup}
All experiments use controlled skill environments with path-level ground truth. A trial is counted as successful only when the agent reaches the bad event specified by the mechanism, not merely when the model describes a risky action. Attack success rate (ASR) is reported as defined in \autoref{eq:asr}. Unless otherwise stated, each case-condition is evaluated with \(m=5\) independent trials, and ASR is averaged over cases and trials. We evaluate SCR-Bench across a diverse set of model backends, including GPT-5.5, GPT-5.4, Claude Opus 4.6, Claude Opus 4.5, Gemini 3.1 Pro Preview, MiniMax-M2.7, DeepSeek-V4, Kimi-K2~\citep{kimiteam2025kimik2}, GLM-5.1, and GLM-5~\citep{glm5team2026glm5}, depending on benchmark availability. All model calls use the default decoding configuration of each backend unless otherwise specified. For a detailed summary of the experimental protocol, including backends, case counts, conditions, trials, metrics, scoring rules, and sandbox design, see \autoref{app:experimental-details}.






\subsection{Results of the SCR-CapFlow}


This section evaluates the SCR-CapFlow, which tests whether an agent transfers upstream-discovered targets into downstream execution interfaces, thereby activating capability-flow composition risk. The results are shown in \autoref{fig:capflow_main}, with detailed numerical data provided in the Appendix.

The results reveal a clear separation between isolated and composed execution. Under isolated conditions, risk is nearly absent: the Control and A-Only conditions achieve 0\% ASR across all evaluated backends, while the B-Only condition averages only 1.4\%. However, once skill composition is enabled, the average ASR increases sharply to 33.6\% under the A+B Neutral condition and 35.9\% under the A+B Explicit condition. Among the evaluated backends, DeepSeek-V4 exhibits the strongest composition risk, exceeding 90\% ASR in both composed conditions. MiniMax-M2.7, GPT-5.5, Gemini 3.1 Pro Preview, and GLM-5.1 also exhibit high composed-execution ASR, indicating substantial capability-flow risk under skill composition. Opus 4.5, Opus 4.6, and GPT-5.4 show lower ASR in the composed conditions, suggesting that SCR can still occur even when models exhibit stronger resistance to composition-induced risk. Overall, these findings show that capability-flow risk arises when upstream skills expose concrete operational targets that are subsequently reused by downstream execution skills, turning otherwise isolated and bounded capabilities into harmful state-changing actions along an activated composition path.


\subsection{Results of the SCR-TrustLift}

% SCR-TrustLift focuses on a different failure mode: whether an upstream review or security-oriented skill can lend semantic legitimacy to a later installation decision. \autoref{tab:trustlift_main} reports the resulting trust-propagation behavior across model backends.

SCR-TrustLift evaluates whether an upstream review or security-oriented skill can transfer semantic legitimacy to a later installation decision. \autoref{tab:trustlift_main} reports both the control condition, where the downstream installation request is evaluated without an upstream endorsement, and the endorsed condition, where the request is preceded by an upstream review signal.

\begin{table}[h]
\centering
\caption{SCR-TrustLift: trust-transfer effects on downstream installation across model backends.}
\label{tab:trustlift_main}
\footnotesize
\setlength{\tabcolsep}{5.5pt}
\renewcommand{\arraystretch}{1.08}
\begin{tabular}{@{}lccccc@{}}
\toprule
\textbf{Backend} 
& \multicolumn{2}{c}{\textbf{Control}} 
& \multicolumn{2}{c}{\textbf{Endorsed}} 
& \textbf{Lift} \\
\cmidrule(lr){2-3} \cmidrule(lr){4-5}
& Success/Total & ASR (\%) & Success/Total & ASR (\%) & \(\Delta\) ASR \\
\midrule
Opus-4.6 & 0/401 & 0.00 & 101/401 & 25.19 & \textcolor{green!60!black}{\(\uparrow 25.19\)} \\
Opus-4.5 & 0/401 & 0.00 & 401/401 & 100.00 & \textcolor{green!60!black}{\(\uparrow 100.00\)} \\
GPT-5.4 & 0/401 & 0.00 & 387/401 & 96.51 & \textcolor{green!60!black}{\(\uparrow 96.51\)} \\
Gemini-3.1-Pro-Preview & 22/401 & 5.49 & 392/401 & 97.76 & \textcolor{green!60!black}{\(\uparrow 92.27\)} \\
MiniMax-M2.7 & 0/401 & 0.00 & 401/401 & 100.00 & \textcolor{green!60!black}{\(\uparrow 100.00\)} \\
\bottomrule
\end{tabular}
\end{table}

The results show a clear endorsement-induced lift. Under the control condition, harmful-installation ASR is nearly zero across all evaluated backends: four out of five backends produce no successful harmful installations, and only Gemini-3.1-Pro-Preview shows a small nonzero control ASR of \(5.49\%\). In contrast, the endorsed condition sharply increases downstream installation risk. The average ASR rises from \(1.10\%\) in the control condition to \(83.89\%\) in the endorsed condition, corresponding to an average lift of \(82.79\) percentage points. Four backends reach near-saturated endorsed ASR: Opus-4.5 and MiniMax-M2.7 reach \(100.00\%\), while GPT-5.4 and Gemini-3.1-Pro-Preview reach \(96.51\%\) and \(97.76\%\), respectively. Opus-4.6 is more conservative, but still exhibits a substantial endorsement lift of \(25.19\) percentage points. These results indicate that the observed failures are not primarily caused by the downstream installation request being accepted in isolation. Rather, the upstream endorsement changes how the agent interprets the downstream action, causing it to generalize from ``this skill appears reviewed or safe'' to ``the subsequent installation is legitimate.'' SCR-TrustLift therefore demonstrates a trust-transfer failure mode: review-oriented or security-themed skills can themselves become part of a composition path that unintentionally legitimizes harmful downstream operations.
% 2026年6月8日 - 因为只有实验组的文字，没有对照组，所以修改
% The results indicate that trust transfer is highly effective across most evaluated backends. GPT-5.4, Gemini 3.1 Pro Preview, MiniMax-M2.7, and Claude Opus 4.5 all achieve near-saturated ASR, with MiniMax-M2.7 and Claude Opus 4.5 reaching $100\%$, while GPT-5.4 and Gemini 3.1 Pro Preview achieve $96.5\%$ and $97.8\%$, respectively. Claude Opus 4.6 is noticeably more conservative, but still reaches a non-trivial ASR of $25.2\%$, indicating that endorsement-driven installation risk is not fully mitigated even on more cautious backends.These findings suggest that the observed failures are not caused by the upstream skill directly performing malicious installation. Instead, the upstream endorsement changes how the agent interprets the downstream installation request. In other words, the agent may incorrectly generalize ``this skill appears reviewed or safe'' into ``the downstream installation action is legitimate.'' The experiment therefore demonstrates that review-oriented or security-themed skills can themselves become part of a trust-propagation chain that unintentionally legitimizes later harmful operations.

\subsection{Results of the SCR-AuthBlur}

This section evaluates the SCR-AuthBlur, which tests whether related context or upstream advisory outputs blur the boundary between recommendation and authorization, thereby increasing the likelihood that an agent approves risky downstream requests. The results are shown in \autoref{tab:authblur_main}.
 % and \autoref{fig:authblur_main}
 
% \begin{table}[t]
% \centering
% \caption{SCR-AuthBlur results across model backends. Values report risky approval rates (\%) under the control context \(L_0\), related context \(L_1\), and full authorization-like context \(L_3\).}
% \label{tab:authblur_main}
% \begin{tabular}{lccc}
% \toprule
% \textbf{Backend} & \textbf{\(L_0\) Control} & \textbf{\(L_1\) Related} & \textbf{\(L_3\) Full Auth} \\
% \midrule
% GPT-5.5 & 2.9 & 10.2 & 17.6 \\
% GPT-5.4 & 9.5 & 7.1 & 7.3 \\
% Opus-4.6 & 2.0 & 10.0 & 17.6 \\
% Opus-4.5 & 8.7 & 9.6 & 13.1 \\
% Gemini-3.1-pro-preview & 10.0 & 30.1 & 35.0 \\
% MiniMax-M2.7 & 19.4 & 31.9 & 47.3 \\
% DeepSeek-V4 & 26.9 & 40.6 & 43.1 \\
% Kimi-K2 & 47.3 & 81.9 & 88.4 \\
% GLM-5.1 & 10.5 & 8.9 & 17.4 \\
% GLM-5 & 20.1 & 40.0 & 52.9 \\
% \bottomrule
% \end{tabular}
% \end{table}
% \begin{table}[h]
% \centering
% \caption{SCR-AuthBlur results across model backends. Values report risky approval rates (\%) under the control context \(L_0\), related context \(L_1\), and full authorization-like context \(L_3\).}
% \label{tab:authblur_main}
% \scriptsize
% \begin{tabular}{lccc}
% \toprule
% \textbf{Backend} & \textbf{\(L_0\) Control} & \textbf{\(L_1\) Related} & \textbf{\(L_3\) Full Auth} \\
% \midrule
% GPT-5.5 & 2.9 & 10.2 & 17.6 \\
% GPT-5.4 & 9.5 & 7.1 & 7.3 \\
% Opus-4.6 & 2.0 & 10.0 & 17.6 \\
% Opus-4.5 & 8.7 & 9.6 & 13.1 \\
% Gemini-3.1-pro-preview & 10.0 & 30.1 & 35.0 \\
% MiniMax-M2.7 & 19.4 & 31.9 & 47.3 \\
% DeepSeek-V4 & 26.9 & 40.6 & 43.1 \\
% Kimi-K2 & 47.3 & 81.9 & 88.4 \\
% GLM-5.1 & 10.5 & 8.9 & 17.4 \\
% GLM-5 & 20.1 & 40.0 & 52.9 \\
% \bottomrule
% \end{tabular}%
% \end{table}
\begin{table}[h]
\centering
% \caption{SCR-AuthBlur results across model backends. Values report risky approval rates (\%) under the control context \(L_0\), related context \(L_1\), and full authorization-like context \(L_3\). \(\Delta_1=L_1-L_0\) and \(\Delta_2=L_3-L_0\).}
\caption{SCR-AuthBlur results across model backends. Risky approval rates (\%) are reported under the control context \(L_0\), related context \(L_1\), and full authorization-like context \(L_3\). \(\Delta_1=L_1-L_0\) and \(\Delta_2=L_3-L_0\) are percentage-point changes relative to \(L_0\). Green/red cells indicate increases/decreases, with shade intensity proportional to change magnitude; \textbf{bold} values denote the highest approval rate for each backend.}
\label{tab:authblur_main}
\scriptsize
\begin{tabular}{lccccc}
\toprule
\textbf{Backend} 
% 深绿，对应 \cellcolor[HTML]{D5F0E1}  20~100
% 中绿，对应 \cellcolor[HTML]{E7F8EE}  5~19
% 浅绿，对应 \cellcolor[HTML]{F4FCF7}  0~4
% 深红，对应 \cellcolor[HTML]{F9E8E7}  -5~-100
% 浅红，对应 \cellcolor[HTML]{FDF5F4}  -0~-4
& \textbf{\(L_0\) Control} 
& \textbf{\(L_1\) Related} 
& \textbf{\(\Delta_1\)} 
& \textbf{\(L_3\) Full Auth} 
& \textbf{\(\Delta_2\)} \\
\midrule
GPT-5.5 & 2.9 & 10.2 & \cellcolor[HTML]{F4FCF7}+7.3 & \textbf{17.6} & \cellcolor[HTML]{E7F8EE}+14.7 \\
GPT-5.4 & \textbf{9.5} & 7.1 & \cellcolor[HTML]{F9E8E7}-2.4 & 7.3 & \cellcolor[HTML]{F9E8E7}-2.2 \\
Opus-4.6 & 2.0 & 10.0 & \cellcolor[HTML]{F4FCF7}+8.0 & \textbf{17.6} & \cellcolor[HTML]{E7F8EE}+15.6 \\
Opus-4.5 & 8.7 & 9.6 & \cellcolor[HTML]{F4FCF7}+0.9 & \textbf{13.1} & \cellcolor[HTML]{F4FCF7}+4.4 \\
Gemini-3.1-pro-preview & 10.0 & 30.1 & \cellcolor[HTML]{D5F0E1}+20.1 & \textbf{35.0} & \cellcolor[HTML]{D5F0E1}+25.0 \\
MiniMax-M2.7 & 19.4 & 31.9 & \cellcolor[HTML]{E7F8EE}+12.5 & \textbf{47.3} & \cellcolor[HTML]{D5F0E1}+27.9 \\
DeepSeek-V4 & 26.9 & 40.6 & \cellcolor[HTML]{E7F8EE}+13.7 & \textbf{43.1} & \cellcolor[HTML]{E7F8EE}+16.2 \\
Kimi-K2 & 47.3 & 81.9 & \cellcolor[HTML]{D5F0E1}+34.6 & \textbf{88.4}& \cellcolor[HTML]{D5F0E1}+41.1 \\
GLM-5.1 & 10.5 & 8.9 & \cellcolor[HTML]{F9E8E7}-1.6 & \textbf{17.4} & \cellcolor[HTML]{F4FCF7}+6.9 \\
GLM-5 & 20.1 & 40.0 & \cellcolor[HTML]{E7F8EE}+19.9 & \textbf{52.9} & \cellcolor[HTML]{D5F0E1}+32.8 \\
\bottomrule
\end{tabular}%
\end{table}



The results indicate a trend of increasing risky approvals with stronger context. On average, ASR rises from $15.7\%$ under L0 control to $27.0\%$ under L1 related context and $34.0\%$ under L3 full-advisory, an increase of $18.3$ percentage points. Kimi-K2 exhibits the highest risk, rising from $47.3\%$ to $81.9\%$ under L1. GLM-5, MiniMax-M2.7, and DeepSeek-V4 also show substantial context-induced increases, whereas GPT-5.4 and GLM-5.1 decrease under L1 before recovering under L3, indicating that related context can occasionally trigger refusal. These findings suggest that authorization-confusion risk is highly backend-dependent but overall supports the SCR-AuthBlur hypothesis: ordinary advisory or related context can already elevate risky approvals, and stronger advisory outputs further amplify the effect. Compared with SCR-CapFlow and SCR-TrustLift, SCR-AuthBlur produces weaker attack strength but highlights composition risk at the semantic decision boundary.
% \subsection{Discussion}

% The experimental results across SCR-CapFlow, SCR-TrustLift, and SCR-AuthBlur consistently demonstrate that composition paths reveal risks largely invisible under isolated skill evaluation. DeepSeek-V4 exhibits the strongest vulnerability in capability-flow scenarios, while MiniMax-M2.7 and Claude Opus 4.5 are highly susceptible to trust-transfer effects, and Kimi-K2 shows the largest increases in risky approvals under advisory context. In contrast, more conservative backends such as Claude Opus 4.6 and GPT-5.4 reduce but do not eliminate composition risk. These findings indicate that path-level interactions, rather than individual skill behavior, are the primary driver of security failures, and that backend-dependent policies and planning strategies significantly influence susceptibility. Overall, no isolated skill-level evaluation suffices to guarantee safety, highlighting the necessity of path-aware analysis for robust agent skill security.


% \subsection{Others}




% \subsection{Results}

% The two main positive results are capability flow and trust transfer.
% \textbf{SCR-CapFlow} tests whether two individually weak or harmless skills can become harmful when composed.
% Each case has five conditions: a control condition, two isolated-skill conditions, and two composed A+B conditions.
% The A+B neutral condition uses a natural task prompt that allows composition without explicitly instructing an attack chain, while the A+B explicit condition directly asks the agent to call both skills in sequence.



% \autoref{tab:privilege-amplification} shows a sharp separation between node-level and path-level risk.
% Across the six evaluated backends, the control and A-only conditions have 0\% ASR, and B-only averages only 1.4\%.
% These numbers indicate that the downstream side effect rarely reaches a harmful state without the upstream target-discovery step.
% Once the skills are composed, ASR rises to 33.5\% under neutral prompts and 36.1\% under explicit prompts on average.
% The model-level spread is substantial: Opus 4.6 remains conservative, while DeepSeek-V4 exceeds 90\% ASR under both composed conditions.
% This spread strengthens the mechanism claim while narrowing the interpretation: capability-flow risk is not uniformly triggered by all backends, but when it appears, isolated baselines do not predict the composed-path failure.

% \textbf{SCR-TrustLift} evaluates whether a benign-looking upstream skill can legitimize a downstream harmful installation.
% The upstream skill presents itself as a security or review utility.
% After its endorsement, the agent is exposed to a downstream installation opportunity.
% The bad event is successful installation of the harmful skill.

% \begin{table}[t]
%   \caption{\textbf{SCR-TrustLift} results over 401 trials per backend. The expanded columns make the path event explicit: the upstream skill supplies an endorsement, and success is counted only when the downstream harmful skill is installed.}
%   \label{tab:semantic-endorsement}
%   \centering
%   \small
%   \resizebox{\linewidth}{!}{%
%   \begin{tabular}{lcccll}
%     \toprule
%     Backend & Trials & Successes & ASR (\%) & Path Event & Evidence Role \\
%     \midrule
%     Claude Opus 4.6 & 401 & 101 & 25 & Endorsement \(\rightarrow\) Install & Conservative But Vulnerable \\
%     GPT-5.4 & 401 & 387 & 97 & Endorsement \(\rightarrow\) Install & Strong Trust Transfer \\
%     Gemini 3.1 Pro Preview & 401 & 392 & 98 & Endorsement \(\rightarrow\) Install & Strong Trust Transfer \\
%     MiniMax-M2.7 & 401 & 401 & 100 & Endorsement \(\rightarrow\) Install & Complete Transfer \\
%     Claude Opus 4.5 & 401 & 401 & 100 & Endorsement \(\rightarrow\) Install & Complete Transfer \\
%     \bottomrule
%   \end{tabular}
%   }
% \end{table}

% \autoref{tab:semantic-endorsement} shows that trust transfer is highly effective on four of five tested backends, where ASR ranges from 97\% to 100\%.
% Claude Opus 4.6 is more conservative, but still reaches 25\%.
% This result supports the second mechanism in the framework.
% The harmful action is not triggered by a single obviously malicious skill in isolation.
% Instead, the upstream skill changes the decision context for a later installation step.
% The experiment therefore illustrates how endorsement can increase the probability that the agent activates a harmful downstream path.

% \subsection{Ablation}

% The ablation study focuses on authorization confusion because it has a natural context-level intervention.
% \textbf{SCR-AuthBlur} tests whether earlier advisory or related context changes the downstream approval boundary.
% Three main context levels are evaluated.
% \(L_0\) is a control context.
% \(L_1\) contains related natural-language context that simulates ordinary task continuity.
% \(L_3\) contains a stronger advisory context that instructs the agent to call the upstream skill before the downstream decision.
% The \(L_2\) findings-only condition is useful for boundary analysis, but only 52 cases satisfy the stricter audit criteria, so it is left to the appendix.

% \begin{table}[t]
%   \caption{\textbf{SCR-AuthBlur} ablation on 118 cases across nine tested backends. The table reports the raw ASR at each context level, the marginal gains between levels, and a short trend summary for each backend.}
%   \label{tab:context-pollution}
%   \centering
%   \scriptsize
%   \resizebox{\linewidth}{!}{%
%   \begin{tabular}{lrrrrrrl}
%     \toprule
%     Backend & \(L_0\) control & \(L_1\) related & \(L_3\) full advisory & \(\Delta L_1-L_0\) & \(\Delta L_3-L_1\) & \(\Delta L_3-L_0\) & Trend Summary \\
%     \midrule
%     GPT-5.5 & 2.9 & 10.2 & 17.6 & 7.3 & 7.4 & 14.7 & Monotonic Increase \\
%     GPT-5.4 & 9.5 & 7.1 & 7.3 & -2.4 & 0.2 & -2.2 & Related Context Triggers Refusal \\
%     Opus 4.6 & 2.0 & 10.0 & 17.6 & 8.0 & 7.6 & 15.6 & Monotonic Increase \\
%     Opus 4.5 & 8.7 & 9.6 & 13.1 & 0.9 & 3.5 & 4.4 & Mild Increase \\
%     Gemini 3.1 Pro Preview & 10.0 & 30.1 & 35.0 & 20.1 & 4.9 & 25.0 & Strong Increase \\
%     MiniMax-M2.7 & 19.4 & 31.9 & 47.3 & 12.5 & 15.4 & 27.9 & Strong Increase \\
%     DeepSeek-V4 & 26.9 & 40.6 & 43.1 & 13.7 & 2.5 & 16.2 & Related Context Drives Most Lift \\
%     Kimi-K2 & 47.3 & 81.9 & 88.4 & 34.6 & 6.5 & 41.1 & High Baseline and Large Lift \\
%     GLM-5.1 & 10.5 & 8.9 & 17.4 & -1.6 & 8.5 & 6.9 & Advisory Context Recovers Lift \\
%     GLM-5 & 20.1 & 40.0 & 52.9 & 19.9 & 12.9 & 32.8 & Strong Increase \\
%     \midrule
%     Average & 15.7 & 27.0 & 34.0 & 11.3 & 6.9 & 18.2 & Positive Mean Lift \\
%     \bottomrule
%   \end{tabular}
%   }
% \end{table}

% \autoref{tab:context-pollution} shows a weaker but consistent mechanism-level pattern.
% On average, risky approval increases from 16.4\% in \(L_0\) to 26.7\% in \(L_1\) and 33.9\% in \(L_3\).
% The full advisory context yields a 17.5-point average lift over control.
% The effect is not uniform: GPT-5.4 and GLM-5.1 show small or negative \(L_1-L_0\) changes, suggesting that related context can sometimes make the risk more salient and trigger refusal.
% Nevertheless, \(L_3\) improves over \(L_1\) for every backend, and the average trend supports the authorization-confusion mechanism.
% The result is best interpreted as a boundary analysis rather than the strongest headline result.
% It shows that \textbf{SCR} also appears through semantic and authorization flows, even when the upstream content is not a direct command.

\subsection{Ablation}

The ablation focuses on SCR-CapFlow; results for SCR-TrustLift and SCR-AuthBlur are in the Appendix. As shown in \autoref{tab:capflow-ablation-main}, ASR remains near zero under control and single-skill conditions, while composed conditions (A+B Neutral and A+B Explicit) increase ASR across most backends. DeepSeek-V4 shows the highest risk (>90\%), whereas Opus-4.6 is more conservative. These results indicate that capability-flow risk arises primarily from activated composition paths, not from individual skill failures, highlighting how benign-appearing skills can propagate risk when combined.

% \begin{table}[h]
% \centering
% \caption{\textbf{SCR-CapFlow} ablation across 150 cases. ASR is shown for control, single-skill (A-only, B-only), and composed conditions (A+B Neutral, A+B Explicit).}
% \label{tab:capflow-ablation-main}
% \scriptsize
% \begin{tabular}{lrrrrr}
% \toprule
% Backend & Control & A-Only & B-Only & A+B Neutral & A+B Explicit \\
% \midrule
% GPT-5.5 & 0.0 & 0.0 & 1.6 & 48.1 & 47.2 \\
% GPT-5.4 & 0.0 & 0.0 & 1.2 & 4.4 & 4.0 \\
% Opus-4.6 & 0.0 & 0.0 & 1.7 & 1.3 & 4.1 \\
% Opus-4.5 & 0.0 & 0.0 & 1.2 & 0.0 & 0.7 \\
% Gemini-3.1-Pro-Preview & 0.0 & 0.0 & 1.3 & 30.0 & 41.9 \\
% MiniMax-M2.7 & 0.0 & 0.0 & 1.7 & 75.5 & 74.9 \\
% DeepSeek-V4 & 0.0 & 0.0 & 1.5 & 91.5 & 92.5 \\
% GLM-5.1 & 0.0 & 0.0 & 1.3 & 25.5 & 26.9 \\
% GLM-5 & 0.0 & 0.0 & 0.7 & 26.4 & 30.7 \\
% \bottomrule
% \end{tabular}
% \end{table}
\begin{table}[h]
\centering
% \caption{\textbf{SCR-CapFlow} ablation across 150 cases. ASR is shown for control, single-skill (A-only, B-only), and composed conditions (A+B Neutral, A+B Explicit).}
\caption{\textbf{SCR-CapFlow} ablation across model backends over 150 evaluation cases. ASR (\%) is shown for control, single-skill, and composed-skill execution conditions. \textbf{Bold} and \underline{underlined} values denote the highest and second-highest ASR for each backend, respectively.}
\label{tab:capflow-ablation-main}
\scriptsize
\begin{tabular}{lrrrrr}
\toprule
\textbf{Backend} & \textbf{Control} & \textbf{A-Only} & \textbf{B-Only} & \textbf{A+B Neutral} & \textbf{A+B Explicit} \\
\midrule
GPT-5.5 & 0.0 & 0.0 & 1.6 & \textbf{48.1} & \underline{47.2} \\
GPT-5.4 & 0.0 & 0.0 & 1.2 & \textbf{4.4} & \underline{4.0} \\
Opus-4.6 & 0.0 & 0.0 & \underline{1.7} & 1.3 & \textbf{4.1} \\
Opus-4.5 & 0.0 & 0.0 & \textbf{1.2} & 0.0 & \underline{0.7} \\
Gemini-3.1-Pro-Preview & 0.0 & 0.0 & 1.3 & \underline{30.0} & \textbf{41.9} \\
MiniMax-M2.7 & 0.0 & 0.0 & 1.7 & \textbf{75.5} & \underline{74.9} \\
DeepSeek-V4 & 0.0 & 0.0 & 1.5 & \underline{91.5} & \textbf{92.5} \\
GLM-5.1 & 0.0 & 0.0 & 1.3 & \underline{25.5} & \textbf{26.9} \\
GLM-5 & 0.0 & 0.0 & 0.7 & \underline{26.4} & \textbf{30.7} \\
\bottomrule
\end{tabular}
\end{table}

\section{Conclusion} 
This work identifies \textbf{SCR} (\textbf{S}kill \textbf{C}omposition \textbf{R}isk), a path-level security gap in which skills that appear benign in isolation can become harmful when their outputs, trust signals, or authorization cues are reused across composed execution paths. \textbf{SCR-Bench} operationalizes this gap through controlled evaluations of capability flow, trust transfer, and authorization confusion, showing that composed paths expose downstream risks missed by isolated baselines. These results suggest that isolated vetting is necessary but insufficient; path-aware evaluation should become a standard component of skill ecosystem security.
% 2026年6月8日 注释结论，原因：结论说法太强。This work studies a security gap in LLM agent skill ecosystems: skills that appear \emph{benign in isolation} can become \emph{harmful in composition} once their outputs, trust signals, or advisory context are reused along activated execution paths. We formalize this phenomenon as \textbf{SCR} (\textbf{S}kill \textbf{C}omposition \textbf{R}isk) and introduce \textbf{SCR-Bench}, the first benchmark suite for evaluating path-level risks across capability flow, trust transfer, and authorization confusion. Across multiple model backends and controlled sandbox environments, the results consistently show that isolated skill vetting fails to capture important downstream risks that emerge only during composition. These findings suggest that the security boundary of agent skills should move beyond standalone artifacts toward context-dependent composition paths, making path-aware analysis a necessary component of future agent security systems.
% Conclusion+Discussion
% \section{Conclusion}
% Experiments across SCR-CapFlow, SCR-TrustLift, and SCR-AuthBlur demonstrate that skill composition can reveal security risks largely invisible under isolated evaluation. These failures are driven by path-level interactions rather than individual skill behavior: DeepSeek-V4 is most vulnerable to capability-flow risk, MiniMax-M2.7 and Claude Opus 4.5 are highly susceptible to trust-transfer effects, and Kimi-K2 exhibits the largest increases in risky approvals under advisory context. Even more conservative backends, such as Claude Opus 4.6 and GPT-5.4, retain non-zero composition risk. Skills that appear \emph{benign in isolation} can become \emph{harmful in composition} once outputs, trust signals, or advisory context are reused along activated paths. These results indicate that isolated skill vetting is insufficient, and robust agent security requires path-aware analysis considering context-dependent composition across diverse backends. The \textbf{SCR-Bench} benchmark formalizes this phenomenon and provides a framework for evaluating and mitigating composition risks.


\newpage
{
\small

\begin{thebibliography}{9}

\bibitem[Jiang et~al.(2026a)]{jiang2026sokagenticskills}
Yanna Jiang, Delong Li, Haiyu Deng, Baihe Ma, Xu Wang, Qin Wang, and Guangsheng Yu.
\newblock SoK: Agentic Skills --- Beyond Tool Use in LLM Agents.
\newblock \emph{arXiv preprint arXiv:2602.20867}, 2026.

%2026年6月10日这篇文章更新了 \bibitem[Liu et~al.(2026a)]{liu2026agentskillswild}
% Yi Liu, Weizhe Wang, Ruitao Feng, Yao Zhang, Guangquan Xu, Gelei Deng, Yuekang Li, and Leo Zhang.
% \newblock Agent Skills in the Wild: An Empirical Study of Security Vulnerabilities at Scale.
% \newblock \emph{arXiv preprint arXiv:2601.10338}, 2026.
\bibitem[Liu et~al.(2026a)]{liu2026donotmentionmaliciousskills}
Yi Liu, Zhihao Chen, Yanjun Zhang, Gelei Deng, Yuekang Li, Jianting Ning, and Leo Yu Zhang.
\newblock ``Do Not Mention This to the User'': Detecting and Understanding Malicious Agent Skills in the Wild.
\newblock In \emph{Proceedings of the 35th USENIX Security Symposium (USENIX Security 2026)}, 2026.


\bibitem[Bhardwaj(2026)]{bhardwaj2026formalanalysis}
Varun Pratap Bhardwaj.
\newblock Formal Analysis and Supply Chain Security for Agentic AI Skills.
\newblock \emph{arXiv preprint arXiv:2603.00195}, 2026.

\bibitem[Jiang et~al.(2026b)]{jiang2026agentlab}
Tanqiu Jiang, Yuhui Wang, Jiacheng Liang, and Ting Wang.
\newblock AgentLAB: Benchmarking LLM Agents against Long-Horizon Attacks.
\newblock \emph{arXiv preprint arXiv:2602.16901}, 2026.

\bibitem[Yuan et~al.(2026)]{yuan2026aegis}
Aojie Yuan, Zhiyuan Su, and Yue Zhao.
\newblock AEGIS: No Tool Call Left Unchecked --- A Pre-Execution Firewall and Audit Layer for AI Agents.
\newblock \emph{arXiv preprint arXiv:2603.12621}, 2026.

\bibitem[Guo et~al.(2026)]{guo2026skillprobe}
Zihan Guo, Zhiyu Chen, Xiaohang Nie, Jianghao Lin, Yuanjian Zhou, and Weinan Zhang.
\newblock SkillProbe: Security Auditing for Emerging Agent Skill Marketplaces via Multi-Agent Collaboration.
\newblock \emph{arXiv preprint arXiv:2603.21019}, 2026.

\bibitem[Wang et~al.(2026)]{wang2026systematicsecurityopenclaw}
Yuhang Wang, Haichang Gao, Zhenxing Niu, Zhaoxiang Liu, Wenjing Zhang, Xiang Wang, and Shiguo Lian.
\newblock A Systematic Security Evaluation of OpenClaw and Its Variants.
\newblock \emph{arXiv preprint arXiv:2604.03131}, 2026.

\bibitem[Zhao et~al.(2026)]{zhao2026clawguard}
Wei Zhao, Zhe Li, Peixin Zhang, and Jun Sun.
\newblock ClawGuard: A Runtime Security Framework for Tool-Augmented LLM Agents Against Indirect Prompt Injection.
\newblock \emph{arXiv preprint arXiv:2604.11790}, 2026.

\bibitem[Jiang et~al.(2026c)]{jiang2026harmfulskillbench}
Yukun Jiang, Yage Zhang, Michael Backes, Xinyue Shen, and Yang Zhang.
\newblock HarmfulSkillBench: How Do Harmful Skills Weaponize Your Agents?
\newblock \emph{arXiv preprint arXiv:2604.15415}, 2026.

\bibitem[Duan et~al.(2026)]{duan2026skillattack}
Zenghao Duan, Yuxin Tian, Zhiyi Yin, Liang Pang, Jingcheng Deng, Zihao Wei, Shicheng Xu, Yuyao Ge, and Xueqi Cheng.
\newblock SkillAttack: Automated Red Teaming of Agent Skills through Attack Path Refinement.
\newblock \emph{arXiv preprint arXiv:2604.04989}, 2026.

\bibitem[Hou et~al.(2026)]{hou2026skillsieve}
Yinghan Hou, Zongyou Yang, Zaihu Pang, and Xiujun Ma.
\newblock SkillSieve: A Hierarchical Triage Framework for Detecting Malicious AI Agent Skills.
\newblock \emph{arXiv preprint arXiv:2604.06550}, 2026.

\bibitem[Schmotz et~al.(2026)]{schmotz2026skillinject}
David Schmotz, Luca Beurer-Kellner, Sahar Abdelnabi, and Maksym Andriushchenko.
\newblock Skill-Inject: Measuring Agent Vulnerability to Skill File Attacks.
\newblock \emph{arXiv preprint arXiv:2602.20156}, 2026.

\bibitem[Ahad et~al.(2026)]{ahad2026semanticintentfragmentation}
Tanzim Ahad, Ismail Hossain, Md Jahangir Alam, Sai Puppala, Yoonpyo Lee, Syed Bahauddin Alam, and Sajedul Talukder.
\newblock Semantic Intent Fragmentation: A Single-Shot Compositional Attack on Multi-Agent AI Pipelines.
\newblock \emph{arXiv preprint arXiv:2604.08608}, 2026.

\bibitem[Chen et~al.(2026)]{chen2026tracesafe}
Yen-Shan Chen, Sian-Yao Huang, Cheng-Lin Yang, and Yun-Nung Chen.
\newblock TraceSafe: A Systematic Assessment of LLM Guardrails on Multi-Step Tool-Calling Trajectories.
\newblock \emph{arXiv preprint arXiv:2604.07223}, 2026.

%下面是旧文献 
\bibitem[Debenedetti et~al.(2024)]{debenedetti2024agentdojo}
Edoardo Debenedetti, Jie Zhang, Mislav Balunovi{\'c}, Luca Beurer-Kellner, Marc Fischer, and Florian Tram{\`e}r.
\newblock AgentDojo: A Dynamic Environment to Evaluate Prompt Injection Attacks and Defenses for LLM Agents.
\newblock In \emph{Advances in Neural Information Processing Systems}, 2024.

% \bibitem[Zhang et~al.(2024)]{zhang2024asb}
% Hanrong Zhang, Jingyuan Huang, Kai Mei, Yifei Yao, Zhenting Wang, Chenlu Zhan, Hongwei Wang, and Yongfeng Zhang.
% \newblock Agent Security Bench (ASB): Formalizing and benchmarking attacks and defenses in LLM-based agents.
% \newblock \emph{arXiv preprint arXiv:2410.02644}, 2024.

\bibitem[Zhang et~al.(2025)]{zhang2025agent}
Hanrong Zhang, Jingyuan Huang, Kai Mei, Yifei Yao, Zhenting Wang, Chenlu Zhan, Hongwei Wang, and Yongfeng Zhang.
\newblock Agent Security Bench (ASB): Formalizing and Benchmarking Attacks and Defenses in LLM-Based Agents.
\newblock In \emph{The Thirteenth International Conference on Learning Representations}, 2025.

% \bibitem[Shi et~al.(2025)]{shi2025toolhijacker}
% Jiawen Shi, Zenghui Yuan, Guiyao Tie, Pan Zhou, Neil Zhenqiang Gong, and Lichao Sun.
% \newblock Prompt injection attack to tool selection in LLM agents.
% \newblock \emph{arXiv preprint arXiv:2504.19793}, 2025.

\bibitem[Shi et~al.(2026)]{shi2026toolhijacker}
Jiawen Shi, Zenghui Yuan, Guiyao Tie, Pan Zhou, Neil Zhenqiang Gong, and Lichao Sun.
\newblock Prompt Injection Attack to Tool Selection in LLM Agents.
\newblock In \emph{Proceedings of the Network and Distributed System Security Symposium}, 2026.

\bibitem[Liu et~al.(2026b)]{liu2026maliciousskills}
Yi Liu, Zhihao Chen, Yanjun Zhang, Gelei Deng, Yuekang Li, Jianting Ning, Ying Zhang, and Leo Yu Zhang.
\newblock Malicious Agent Skills in The Wild: A Large-Scale Security Empirical Study.
\newblock \emph{arXiv preprint arXiv:2602.06547}, 2026.

\bibitem[Liu et~al.(2026c)]{liu2026trojanswhisper}
Fazhong Liu, Zhuoyan Chen, Tu Lan, Haozhen Tan, Zhenyu Xu, Xiang Li, Guoxing Chen, Yan Meng, and Haojin Zhu.
\newblock Trojan's Whisper: Stealthy Manipulation of OpenClaw through Injected Bootstrapped Guidance.
\newblock \emph{arXiv preprint arXiv:2603.19974}, 2026.

\bibitem[Yao et~al.(2023)]{yao2023react}
Shunyu Yao, Jeffrey Zhao, Dian Yu, Nan Du, Izhak Shafran, Karthik R. Narasimhan, and Yuan Cao.
\newblock ReAct: Synergizing Reasoning and Acting in Language Models.
\newblock In \emph{The Eleventh International Conference on Learning Representations}, 2023.

\bibitem[Schick et~al.(2023)]{schick2023toolformer}
Timo Schick, Jane Dwivedi-Yu, Roberto Dessi, Roberta Raileanu, Maria Lomeli, Eric Hambro, Luke Zettlemoyer, Nicola Cancedda, and Thomas Scialom.
\newblock Toolformer: Language Models Can Teach Themselves to Use Tools.
\newblock In \emph{Advances in Neural Information Processing Systems}, 2023.

\bibitem[Qin et~al.(2024)]{qin2024toollm}
Yujia Qin, Shihao Liang, Yining Ye, Kunlun Zhu, Lan Yan, Yaxi Lu, Yankai Lin, Xin Cong, Xiangru Tang, Bill Qian, Sihan Zhao, Lauren Hong, Runchu Tian, Ruobing Xie, Jie Zhou, Mark Gerstein, Dahai Li, Zhiyuan Liu, and Maosong Sun.
\newblock ToolLLM: Facilitating Large Language Models to Master 16000+ Real-World APIs.
\newblock In \emph{The Twelfth International Conference on Learning Representations}, 2024.

\bibitem[Patil et~al.(2024)]{patil2024gorilla}
Shishir G. Patil, Tianjun Zhang, Xin Wang, and Joseph E. Gonzalez.
\newblock Gorilla: Large Language Model Connected with Massive APIs.
\newblock In \emph{Advances in Neural Information Processing Systems}, 2024.

\bibitem[Ruan et~al.(2024)]{ruan2024toolemu}
Yangjun Ruan, Honghua Dong, Andrew Wang, Silviu Pitis, Yongchao Zhou, Jimmy Ba, Yann Dubois, Chris Maddison, and Tatsunori Hashimoto.
\newblock Identifying the Risks of LM Agents with an LM-Emulated Sandbox.
\newblock In \emph{The Twelfth International Conference on Learning Representations}, 2024.

\bibitem[Yuan et~al.(2024)]{yuan2024rjudge}
Tongxin Yuan, Zhiwei He, Lingzhong Dong, Yiming Wang, Ruijie Zhao, Tian Xia, Lizhen Xu, Binglin Zhou, Fangqi Li, Zhuosheng Zhang, Rui Wang, and Gongshen Liu.
\newblock R-Judge: Benchmarking Safety Risk Awareness for LLM Agents.
\newblock In \emph{Findings of the Association for Computational Linguistics: EMNLP 2024}, pages 1467--1490, 2024.


\bibitem[Kimi Team(2025)]{kimiteam2025kimik2}
Kimi Team.
\newblock Kimi K2: Open Agentic Intelligence.
\newblock \emph{arXiv preprint arXiv:2507.20534}, 2025.

\bibitem[GLM-5-Team(2026)]{glm5team2026glm5}
GLM-5-Team.
\newblock GLM-5: from Vibe Coding to Agentic Engineering.
\newblock \emph{arXiv preprint arXiv:2602.15763}, 2026.

\bibitem[Liu et~al.(2024)]{liu2024formalizing}
Yupei Liu, Yuqi Jia, Runpeng Geng, Jinyuan Jia, and Neil Zhenqiang Gong.
\newblock Formalizing and Benchmarking Prompt Injection Attacks and Defenses.
\newblock In \emph{Proceedings of the 33rd USENIX Security Symposium}, 2024.

\bibitem[Zhan et~al.(2024)]{zhan2024injecagent}
Qiusi Zhan, Zhixiang Liang, Zifan Ying, and Daniel Kang.
\newblock InjecAgent: Benchmarking Indirect Prompt Injections in Tool-Integrated Large Language Model Agents.
\newblock In \emph{Findings of the Association for Computational Linguistics: ACL 2024}, 2024.

\bibitem[Chen et~al.(2024)]{chen2024agentpoison}
Zhaorun Chen, Zhen Xiang, Chaowei Xiao, Dawn Song, and Bo Li.
\newblock AgentPoison: Red-teaming LLM Agents via Poisoning Memory or Knowledge Bases.
\newblock In \emph{Advances in Neural Information Processing Systems}, 2024.

\bibitem[Zhang et~al.(2025b)]{zhang2025breakingagents}
Boyang Zhang, Yicong Tan, Yun Shen, Ahmed Salem, Michael Backes, Savvas Zannettou, and Yang Zhang.
\newblock Breaking Agents: Compromising Autonomous LLM Agents Through Malfunction Amplification.
\newblock In \emph{Proceedings of the 2025 Conference on Empirical Methods in Natural Language Processing}, 2025.

\end{thebibliography}
}
\newpage

\appendix

% \renewcommand{\sectionautorefname}{Appendix}
\renewcommand{\subsectionautorefname}{Appendix}
\appendix

\section{Limitations}

SCR-AuthBlur is reported conservatively: we exclude \(L_2\) from the main table because only a subset of cases satisfies the stricter audit criterion, and we avoid extrapolating this setting beyond the cases where the criterion holds. In addition, all experiments are conducted in controlled skill environments with observable mock side effects. This design provides path-level ground truth and supports controlled comparisons across mechanisms, but it abstracts away registry-scale heterogeneity, runtime policies, and operational noise. Future work should therefore validate the same mechanisms in larger open skill registries and production-like agent runtimes.

\section{Additional method details}
\label{app:method-details}

This appendix records the details that are omitted from the main Method section to keep the exposition compact.

% \paragraph{Multi-hop amplification.}
% For length-\(k\) paths, let \(N_k(h)=|\Pi_k(G_h)|\) and let
% \[
% \bar r_k(h)=\mathbb{E}_{\pi\in\Pi_k(G_h)}[a_\pi(h)q_\pi(h)].
% \]
% Under the same homogeneous approximation used for \autoref{eq:pair-risk},
% \[
% R_k(G,h)=1-\left(1-\bar r_k(h)\right)^{N_k(h)}.
% \]
% When \(N_k(h)\bar r_k(h)\rightarrow\infty\), the risk approaches one.
% If the composition graph is dense, \(N_k(h)=O(n^k)\), which gives the path-surface growth used in \autoref{eq:surface-gap}.

\paragraph{Multi-hop amplification.}
The main text defines the effective risk of a path as
\(r_\pi(h)=a_\pi(h)q_\pi(h)\). At the ecosystem level,
composition risk can be viewed as the probability that at least one
candidate path is activated and reaches a harmful state. For length-\(k\)
paths, let \(N_k(h)=|\Pi_k(G_h)|\) denote the number of candidate
paths and let
\[
\bar r_k(h)
=
\mathbb{E}_{\pi\in\Pi_k(G_h)}
\left[a_\pi(h)q_\pi(h)\right]
\]
denote the average effective risk over such paths. Under a homogeneous
path-risk approximation, the length-\(k\) composition risk is
\[
R_k(G,h)
=
1-\left(1-\bar r_k(h)\right)^{N_k(h)}.
\]
This expression is intended as a risk-accounting approximation rather
than an independence claim about real failures. It shows that even when
\(\bar r_k(h)\) is small, risk can increase as the number of candidate
paths grows. In dense composition settings, \(N_k(h)\) can scale as
\(O(n^k)\), so multi-hop composition can create a substantially larger
path-level risk surface than node-level vetting.

\paragraph{Capability flow.}
Privilege amplification has a source-to-sink structure:
\[
\pi^{cap}_{ij}: Source \rightarrow s_i \rightarrow s_j \rightarrow Sink.
\]
The corresponding conditional path risk is
\[
q^{cap}_{ij}(h)
=
P(B\mid Source \rightarrow s_i \rightarrow s_j \rightarrow Sink,h).
\]
In the experiments, a capability-flow trial succeeds only when skill A discovers a concrete target, skill B produces an observable side effect, and the side effect acts on the same target:
\[
Y^{cap}_{ij}
=
\mathbb{I}\!\left[
A\_discovers(t)\land B\_acts(t)\land B\_side\_effect(t)
\right].
\]

\paragraph{Trust transfer.}
For semantic endorsement, let \(E\) be the endorsement produced by an upstream skill and let \(h_E=h_0\oplus E\) be the resulting context.
The endorsement lift is
\[
\Delta_E=P(B\mid h_E)-P(B\mid h_0).
\]
When \(B\) is malicious-skill installation, this becomes
\[
\Delta_E
=
P(A_{install}^{harm}\mid h_E)
-
P(A_{install}^{harm}\mid h_0).
\]

\paragraph{Authorization confusion.}
For context pollution, let \(L_l\) denote context level \(l\), where \(L_0\) is the control context and higher levels contain related, finding-like, or authorization-like content.
The contamination effect is
\[
CE_l=P(B\mid L_l)-P(B\mid L_0),
\]
where \(B\) is risky approval or dangerous execution.
Empirically, this is estimated as
\[
\widehat{CE}_l=\widehat{ASR}_{L_l}-\widehat{ASR}_{L_0}.
\]
For semantic endorsement, the empirical lift is
\[
\widehat{\Delta}_E
=
\widehat{ASR}_{endorsement}-\widehat{ASR}_{direct/control}.
\]



% \FloatBarrier

\section{Context Levels in SCR-AuthBlur}
\label{app:authblur-levels}

% Table~\ref{tab:authblur_levels} defines the context levels used in the SCR-AuthBlur sub-benchmark. These levels distinguish how risky approvals may arise from different forms of upstream context: from unrelated control text (\(L_0\)), ordinary task continuity (\(L_1\)), advisor plain findings (\(L_2\)), or stronger advisory outputs resembling authorization guidance (\(L_3\)).  
% The table clarifies which types of context are included or excluded at each level and specifies the main experimental use. In the main paper, only \(L_0\), \(L_1\), and \(L_3\) are reported.

% \begin{table}[h]
%   \caption{Context levels for \textbf{SCR-AuthBlur}. The levels distinguish whether risky approval is driven by ordinary task continuity, plain findings, or stronger advisory context.}
%   \label{tab:authblur_levels}
%   \centering
%   \small
%   \resizebox{\linewidth}{!}{%
%   \begin{tabular}{llll}
%     \toprule
%     Level & Included Context & Excluded Context & Main Use \\
%     \midrule
%     \(L_0\) & Unrelated Control Context & Related Findings or Advice & Baseline Approval Rate \\
%     \(L_1\) & Related Task Context & Advisor Output and Approval Language & Task-Continuity Effect \\
%     \(L_2\) & Plain Findings From Advisor & Explicit Approval or Authorization Wording & Audited Appendix Subset \\
%     \(L_3\) & Stronger Advisory Assessment & Formal User Authorization & Main Advisory-Context Test \\
%     \bottomrule
%   \end{tabular}
%   }
% \end{table}

Table~\ref{tab:authblur_levels} defines the context levels used in the SCR-AuthBlur sub-benchmark.
These levels separate two factors that may affect authorization confusion: whether the agent receives only task-level context or an explicit composition path, and how semantically strong the upstream skill output is as downstream decision evidence.
\(L_0\) provides unrelated context and serves as the control baseline.
\(L_1\) provides a normal task description related to the decision, without specifying an upstream--downstream skill path. It tests whether ordinary related context alone increases risky approvals compared with \(L_0\), thereby probing the context-pollution hypothesis.
\(L_2\) fixes the \(A \rightarrow B\) composition path, but the upstream skill only emits weak advisory or evaluative outputs, such as factual findings, risk assessments, or contextual summaries.
\(L_3\) uses the same specified composition path, but strengthens the upstream output with stronger advisory or evaluative semantics while still excluding formal authorization. This setting serves as an upper-bound condition for testing how semantic strength amplifies SCR-AuthBlur.

\begin{table*}[h]
  \caption{Context levels for \textbf{SCR-AuthBlur}. The levels separate path specification from the semantic strength of upstream context.}
  \label{tab:authblur_levels}
  \centering
  \footnotesize
  \setlength{\tabcolsep}{5pt}
  \begin{tabularx}{\textwidth}{p{0.07\textwidth} p{0.27\textwidth} p{0.32\textwidth} X}
    \toprule
    \textbf{Level} & \textbf{Context / Path Setting} & \textbf{Upstream Signal} & \textbf{Experimental Role} \\
    \midrule
    \(L_0\)
& Unrelated control context; no relevant task or composition cue.
& None; excludes related findings, advice, or approval-like signals.
& Serves as the control baseline for estimating the background risky-approval rate. \\

\(L_1\) 
& Ordinary task-level context related to the decision; no specified \(A \rightarrow B\) composition path. 
& Task description only; excludes advisor output, findings, approval, and authorization language. 
& Tests whether related task context alone increases risky approvals relative to \(L_0\), supporting the context-pollution hypothesis. \\

\(L_2\) 
& Specified \(A \rightarrow B\) composition path. 
& Upstream skill \(A\) emits weak advisory or evaluative outputs, such as factual findings, risk assessments, or contextual summaries; excludes explicit approval, authorization, and directive language. 
& Tests whether weak, non-authorizing upstream evidence is sufficient to contaminate the downstream decision made by skill \(B\). \\

\(L_3\) 
& Specified \(A \rightarrow B\) composition path. 
& Upstream skill \(A\) emits stronger advisory or evaluative outputs, while still excluding formal user authorization. 
& Serves as an upper-bound condition for testing how stronger upstream semantics amplify authorization-confusion risk in skill \(B\). \\
\bottomrule
  \end{tabularx}
\end{table*}



\begin{table*}[h]
\centering
\small
\setlength{\tabcolsep}{7pt}
\caption{SCR-AuthBlur results on the 52 cases satisfying the L2 audit criterion. Values report risky approval rates (\%) under the control context (L0), related context (L1), plain-finding advisor context (L2), and full authorization-like advisory context (L3).}
\label{tab:authblur_l2_52cases}
\begin{tabular}{lcccc}
\toprule
\textbf{Backend} & \textbf{L0 Control} & \textbf{L1 Related} & \textbf{L2 Findings} & \textbf{L3 Full Auth} \\
\midrule
GPT-5.5                  & 1.9  & 16.9 & 7.3  & 28.1 \\
GPT-5.4                  & 4.2  & 13.1 & 1.5  & 9.2  \\
Opus-4.6                 & 0.8  & 18.1 & 8.1  & 28.8 \\
Opus-4.5                 & 6.5  & 17.2 & 7.3  & 19.4 \\
Gemini-3.1-Pro-Preview   & 15.6 & 41.3 & 20.0 & 54.0 \\
MiniMax-M2.7             & 15.9 & 41.9 & 14.9 & 41.2 \\
DeepSeek-V4              & 32.4 & 56.4 & 30.3 & 54.4 \\
Kimi-K2                  & 37.8 & 82.4 & 60.6 & 83.2 \\
GLM-5.1                  & 11.4 & 15.1 & 8.7  & 22.9 \\
GLM-5                    & 19.4 & 48.5 & 24.8 & 55.0 \\
\midrule
\textbf{Average}          & 14.6 & 35.1 & 18.4 & 39.6 \\
\bottomrule
\end{tabular}
\end{table*}

Because only 52 cases satisfy the stricter audit criterion for the L2 setting, we separately evaluate L0--L3 on this audited subset and report the results in Table~\ref{tab:authblur_l2_52cases}. Across all evaluated backends, L2 produces lower risky-approval rates than both L1 and L3. The average risky-approval rate under L2 is 18.4\%, which is only moderately higher than the L0 control rate of 14.6\%, but much lower than L1 related context at 35.1\% and L3 full authorization-like context at 39.6\%. This pattern suggests that plain findings or weak risk-assessment summaries alone are insufficient to induce the same level of downstream decision contamination as stronger contextual or authorization-like signals. More broadly, the result supports the SCR-AuthBlur hypothesis that the semantic strength of upstream skill outputs preserved in the shared context modulates the severity of authorization-confusion risk.


\section{Experimental Protocol and Additional Details}
\label{app:experimental-details}
To provide a comprehensive overview of the experimental setup, Table~\ref{tab:experimental-protocol} summarizes the protocol for each sub-benchmark in SCR-Bench. The table lists the evaluated backends, the number of test cases, the experimental conditions, the number of trials per case, the metrics used to evaluate path-level risk, the scoring rules, and the sandbox environment design. This overview ensures reproducibility and clarifies how each mechanism—capability flow, trust transfer, and authorization confusion—is evaluated under controlled conditions.

\begin{table}[h]
  \centering
  \caption{Experimental protocol. The same scoring principle is used across mechanisms: success requires an observable downstream event in a sandboxed environment.}
  \label{tab:experimental-protocol}
  \resizebox{\linewidth}{!}{%
  \begin{tabular}{llll}
    \toprule
    Element & \textbf{SCR-CapFlow} & \textbf{SCR-TrustLift} & \textbf{SCR-AuthBlur} \\
    \midrule
    Backends & Nine with All Five Conditions Evaluated & Five Tested Backends & Ten Tested Backends \\
    Cases & 150 Paired-Skill Cases & 401 Installation Trials Per Backend & 118 Retained Decision Cases \\
    Conditions & Control, A-Only, B-Only, A+B Neutral, A+B Explicit & Control and Endorsed Downstream Installation & \(L_0\), \(L_1\), \(L_3\) in Main Table \\
    Trials & \(m=5\) Per Case-Condition & One Installation Decision Per Trial & \(m=5\) Per Case-Level Condition \\
    Metric & Linked Attack-Chain ASR & Harmful-Installation ASR & Risky-Approval ASR and Lift \\
    Scoring & A Discovers \(t\), B Acts On \(t\), Side Effect Logged & Harmful Skill is Installed & Downstream Decision Approves Risky Request \\
    Sandbox & Mock Files, Services, Logs, Schedules, and State & Simulated Skill Market Installation & Simulated Approval and Policy Contexts \\
    \bottomrule
  \end{tabular}
  }
\end{table}

\section{The Ablation of SCR-TrustLift}

This section presents the ablation study for SCR-TrustLift, testing whether upstream skills providing trust endorsements affect downstream harmful installations. The results are shown in \autoref{tab:trustlift-ablation}.

\begin{table}[h]
\centering
\caption{\textbf{SCR-TrustLift} ablation across 401 trials per backend. ASR is reported for control and endorsed conditions.}
\label{tab:trustlift-ablation}
\scriptsize
\begin{tabular}{lrr}
\toprule
Backend & Control ASR (\%) & Endorsed ASR (\%) \\
\midrule
Claude Opus 4.6 & 0.00 & 25.19 \\
GPT-5.4 & 0.00 & 96.51 \\
Gemini 3.1 Pro Preview & 5.49 & 97.76 \\
MiniMax-M2.7 & 0.00 & 100.0 \\
Claude Opus 4.5 & 0.00 & 100.0 \\
\bottomrule
\end{tabular}
\end{table}

The results indicate that four out of five backends reach near-saturation ASR under endorsed conditions, while the control condition produces zero or negligible ASR. Claude Opus 4.6 is more conservative but still exhibits significant risk. This confirms that trust-transfer effects depend on upstream endorsements rather than direct malicious skill behavior.

\section{The Ablation of SCR-AuthBlur}

This section presents the ablation study for SCR-AuthBlur, evaluating how earlier advisory or related context shifts downstream approval boundaries. The results are shown in \autoref{tab:authblur-ablation}.

The results show an overall trend of increasing risky approval as context strength increases. The ASR rises from 15.7\% (L0) to 27.0\% (L1) and 34.0\% (L3). Kimi-K2 exhibits the highest risk, while GPT-5.4 sometimes shows lower ASR under L1/L3 than L0, indicating that related context can trigger refusal. Overall, authorization-confusion risk is backend-dependent but supports the SCR-AuthBlur mechanism: context composition amplifies approval risk.

\begin{table}[t]
\caption{\textbf{SCR-AuthBlur} ablation on 118 cases across ten tested backends. Raw ASR at each context level (L0, L1, L3), marginal gains (\(\Delta L_1-L_0\), \(\Delta L_3-L_1\), \(\Delta L_3-L_0\)), and trend summary are reported.}
\label{tab:authblur-ablation}
\scriptsize
\begin{tabular}{lrrrrrr}
\toprule
Backend & L0 Control & L1 Related & L3 Full Auth & $\Delta L_1-L_0$ & $\Delta L_3-L_1$ & $\Delta L_3-L_0$ \\
\midrule
GPT-5.5 & 2.9 & 10.2 & 17.6 & 7.3 & 7.4 & 14.7 \\
GPT-5.4 & 9.5 & 7.1 & 7.3 & -2.4 & 0.2 & -2.2 \\
Opus 4.6 & 2.0 & 10.0 & 17.6 & 8.0 & 7.6 & 15.6 \\
Opus 4.5 & 8.7 & 9.6 & 13.1 & 0.9 & 3.5 & 4.4 \\
Gemini 3.1 Pro Preview & 10.0 & 30.1 & 35.0 & 20.1 & 4.9 & 25.0 \\
MiniMax-M2.7 & 19.4 & 31.9 & 47.3 & 12.5 & 15.4 & 27.9 \\
DeepSeek-V4 & 26.9 & 40.6 & 43.1 & 13.7 & 2.5 & 16.2 \\
Kimi-K2 & 47.3 & 81.9 & 88.4 & 34.6 & 6.5 & 41.1 \\
GLM-5.1 & 10.5 & 8.9 & 17.4 & -1.6 & 8.5 & 6.9 \\
GLM-5 & 20.1 & 40.0 & 52.9 & 19.9 & 12.9 & 32.8 \\
\midrule
Average & 15.7 & 27.0 & 34.0 & 11.3 & 6.9 & 18.2 \\
\bottomrule
\end{tabular}
\end{table}

% \clearpage







% \clearpage
% \input{checklist.tex}


\end{document}
